%==============================================================================
%  sEMG-driven force estimation: learned compensation of an EDM base model
%
%  Compile:  pdflatex main && bibtex main && pdflatex main && pdflatex main
%
%  Class: IEEEtran (IEEE TNSRE / TBME style). To switch to a generic article,
%  replace the \documentclass line with:
%      \documentclass[11pt]{article}
%      \usepackage[margin=1in]{geometry}
%  and delete \IEEEpeerreviewmaketitle below.
%
%  Figures: copy the PNGs you want from
%      mycode/Estimation_2026_10_13/paper_runs/<stamp>/E*/...
%  into paper/figures/ and reference them by bare filename.
%==============================================================================
\documentclass[journal]{IEEEtran}

\usepackage{graphicx}
\usepackage{amsmath}
\usepackage{amssymb}
\usepackage{booktabs}
\usepackage{array}
\usepackage{url}
\usepackage{cite}
\usepackage{siunitx}

\graphicspath{{figures/}}

\begin{document}

\title{Learned Compensation of Electromechanical Delay Model Error for
Cross-Day sEMG-Driven Force Estimation}

\author{Nuwan~Madhusanka,~\IEEEmembership{}
        and~Co-authors%
\thanks{Manuscript received ...; revised ....}%
\thanks{The authors are with La Trobe University, Melbourne, Australia
(e-mail: ...).}}

\maketitle

%==============================================================================
\begin{abstract}
% 150-250 words. State, in order:
%   (i)   problem: EDM models leave systematic residuals in sEMG-force estimation
%   (ii)  method: per-session gain calibration + learned residual compensation
%   (iii) evaluation: leave-one-day-out across 4 sessions, 7 architectures
%   (iv)  headline numbers: 20.5% cross-day RMSE reduction (0.5581 -> 0.4436)
%   (v)   the negative result: architectures differ by only 4.7%
%   (vi)  the boundary: transfer fails outside the training gain range
\end{abstract}

\begin{IEEEkeywords}
Surface electromyography, force estimation, electromechanical delay,
cross-session generalisation, deep learning, state-space models,
physics-informed neural networks.
\end{IEEEkeywords}

\IEEEpeerreviewmaketitle

%==============================================================================
\section{Introduction}
\label{sec:intro}
% Motivation: sEMG-driven force estimation for prosthetics / rehabilitation.
% The EDM (electromechanical delay) base model and its residual error.
% Why cross-day matters: in deployment the model is trained one day and used
% on another.

% State contributions explicitly as a list the reader can check against results.
The contributions of this work are:
\begin{enumerate}
  \item A per-session gain calibration procedure that makes cross-day
        compensation viable, using a short reference contraction available at
        deployment time.
  \item A learned residual compensator that reduces cross-day force estimation
        error by 20.5\% relative to a calibrated EDM baseline.
  \item A controlled comparison of seven sequence architectures showing they
        perform equivalently, indicating that architecture is not the limiting
        factor for this task.
  \item Characterisation of the transfer boundary: compensation degrades when
        the test session's gain falls outside the range spanned by training
        sessions.
\end{enumerate}

%==============================================================================
\section{Related Work}
\label{sec:related}
\subsection{sEMG-driven force estimation}
% Hill-type models, EDM models, regression approaches.

\subsection{Deep sequence models for biosignals}
% CNN/LSTM/Transformer for EMG; state-space models; physics-informed networks.

\subsection{Cross-session generalisation in sEMG}
% Electrode shift, impedance change, day-to-day variability, calibration
% protocols, domain adaptation.

%==============================================================================
%==============================================================================
% Section III: Methodology
%
% Included from main.tex via %==============================================================================
% Section III: Methodology
%
% Included from main.tex via %==============================================================================
% Section III: Methodology
%
% Included from main.tex via \input{sections/methodology}
% Compile main.tex, not this file.
%==============================================================================

%==============================================================================
\section{Methodology}
\label{sec:method}

%------------------------------------------------------------------------------
\subsection{Problem Formulation}
\label{sec:formulation}

Let $x(t)$ denote the preprocessed surface electromyography (sEMG) envelope and
let $y(t)$ denote the simultaneously measured force. A physics-based base model
produces an estimate $\hat{y}_{\mathrm{base}}(t)$ from the sEMG signal alone.
This estimate captures the behaviour that the assumed physiology predicts, but
it leaves a systematic residual. We therefore write the final force estimate as
the sum of the base model output and a learned correction,
\begin{equation}
  \hat{y}(t) \;=\; \hat{y}_{\mathrm{base}}(t) \;+\; \Delta(t),
  \label{eq:decomposition}
\end{equation}
where $\Delta(t)$ is produced by a neural network from a window of recent
observations.

This decomposition is deliberate. The base model carries the part of the
relationship that is known in advance and that transfers between recording
sessions, while the network is asked only to explain what the base model cannot.
Compared with training a network to predict force directly, this keeps the
learned component small, makes its contribution measurable against an explicit
baseline, and ensures that a failure of the network degrades the estimate
towards the physics-based prediction rather than towards an arbitrary one.

%------------------------------------------------------------------------------
\subsection{Dataset}
\label{sec:dataset}

The dataset comprises 42 recordings collected across five sessions. Each
recording contains a time channel, a rectified sEMG channel, a low-pass filtered
sEMG channel, and a measured force channel, sampled at \SI{2}{\kilo\hertz}
($T_s = \SI{0.5}{\milli\second}$). This work uses the low-pass filtered sEMG
channel as the model input and the force channel as the target. Recordings in
four of the sessions are \SI{25}{\second} long (50{,}000 samples) and those in
the remaining session are \SI{15}{\second} long (30{,}000 samples).

Throughout this paper a session refers to a single electrode setup rather than
simply a calendar date. Within one session the electrodes were applied once and
left in place, so all recordings in that session share the same electrode
position, skin preparation and contact impedance. Between sessions the entire
setup was renewed. This distinction matters because the relationship between
sEMG amplitude and produced force is sensitive to exactly those factors.

Table~\ref{tab:sessions} summarises the sessions. For each recording we
characterise the input to output relationship by a single gain, defined as the
ratio of the 99th percentile of force to the 99th percentile of normalised sEMG.
The per-session median of this quantity varies from 0.44 to 5.86, a factor of
roughly thirteen between the extreme sessions, and the ordering is close to
monotonic across the five sessions. Variation of this magnitude means that a
model fitted on one session will systematically over- or under-predict on
another unless the gain is explicitly accounted for. Section~\ref{sec:calibration}
describes how we handle it.

\begin{table}[!t]
\caption{Recording sessions. The gain is the ratio of the 99th percentile force
to the 99th percentile normalised sEMG, reported as the per-session median.}
\label{tab:sessions}
\centering
\begin{tabular}{lrrrrl}
\toprule
Session & Files & Dur.~(s) & $F_{99}$ & Gain & Role \\
\midrule
1st\_Comp  & 5  & 25 & 3.68 & 3.70 & test / val / train \\
2nd\_Comp  & 9  & 25 & 5.56 & 5.86 & test / val / train \\
3rd\_Base  & 5  & 25 & 0.82 & 0.83 & test / val / train \\
4th\_Valid & 13 & 15 & 1.11 & 1.12 & test / val / train \\
5th\_Comp  & 10 & 25 & 0.44 & 0.44 & training only \\
\bottomrule
\end{tabular}
\end{table}

%------------------------------------------------------------------------------
\subsection{Per-Recording sEMG Normalisation}
\label{sec:semgnorm}

Each recording is normalised by its own amplitude statistics before entering the
model. We divide by the 99th percentile of the positive sEMG samples in that
recording,
\begin{equation}
  x_n(t) \;=\; \frac{x(t)}
                    {\mathrm{p}_{99}\{\,x(\tau) \;:\; x(\tau) > 0\,\}} ,
  \label{eq:semgnorm}
\end{equation}
using a percentile rather than the maximum so that isolated spikes do not set
the scale.

This step removes the component of between-session variation that originates in
electrode gain and contact quality. It is important to note that the
normalisation constant is computed from the input signal only. The sEMG signal
is available at deployment time, so the same quantity can be computed for a new
recording without any knowledge of the force that is being predicted. The
normalisation therefore introduces no label leakage, and this remains true when
the recording comes from a session the model has never seen.

%------------------------------------------------------------------------------
\subsection{Base Model}
\label{sec:basemodel}

The base model combines a static saturating nonlinearity with first-order
dynamics. The nonlinearity represents the compressive relationship between
muscle activation and produced force at higher activation levels, and the
dynamics represent the delay and smoothing introduced by the electromechanical
coupling between electrical activity and measurable force. The two stages are
\begin{align}
  z(t) &= g \,\tanh\!\big(s\, x_n(t)\big), \label{eq:tanh} \\
  G(\mathrm{s}) &= \frac{\omega_c}{\mathrm{s} + \omega_c},
                   \qquad \omega_c = 2\pi f_c, \label{eq:lpf}
\end{align}
where $g$ is a gain, $s$ shapes the saturation, and $f_c$ is the cutoff
frequency of the first-order stage. The continuous-time filter in
\eqref{eq:lpf} is discretised with the bilinear transform. Writing
$K = 2/T_s$, the resulting difference equation has numerator coefficients
$b_0 = b_1 = \omega_c/(K + \omega_c)$ and denominator coefficient
$a_1 = (\omega_c - K)/(K + \omega_c)$, and the base model output is
$\hat{y}_{\mathrm{base}} = G\{z\}$.

Including the first-order stage is necessary rather than cosmetic. When force is
plotted against sEMG for a single recording, the trajectory forms a hysteresis
loop: the force observed at a given sEMG level differs depending on whether the
contraction is building or releasing (Fig.~\ref{fig:hysteresis}). A static
nonlinearity is single-valued by construction and cannot reproduce this
behaviour, no matter how its parameters are chosen. The ablation in
Section~\ref{sec:ablation-lpf} quantifies what is lost when the stage is removed.

\begin{figure}[!t]
\centering
% \includegraphics[width=\columnwidth]{force_vs_semg_hysteresis.png}
\caption{Force plotted against normalised sEMG for representative recordings.
The trajectories form hysteresis loops, so the force observed at a given sEMG
level differs between the rising and falling phases of a contraction. A static
nonlinearity is single-valued and cannot represent this, which motivates the
first-order stage in \eqref{eq:lpf}.}
\label{fig:hysteresis}
\end{figure}

The three parameters $\{g, s, f_c\}$ are estimated by nonlinear least squares
using a trust region reflective algorithm, with $g$ and $s$ constrained to
$[0, 50]$ and $f_c$ constrained to $[1, 200]\,\si{\hertz}$. The fit is performed
on the concatenated residuals of all training recordings of a given fold. No
recording that will later be used for validation or testing contributes to the
fit, and the base model is refitted independently for each fold so that this
holds throughout.

%------------------------------------------------------------------------------
\subsection{Per-Session Gain Calibration}
\label{sec:calibration}

Because the input to output gain varies substantially across sessions, a single
fitted base model cannot serve all of them. We therefore estimate one scalar
gain per recording from a short calibration segment, which mirrors the practical
protocol of performing a brief reference contraction when electrodes are fitted.

Let $\mathcal{C}$ denote the set of samples in the calibration segment. The
calibration gain is the least squares scaling that best aligns the base model
output with the measured force over that segment,
\begin{equation}
  \gamma \;=\; \frac{\sum_{t \in \mathcal{C}} y(t)\, \hat{y}_{\mathrm{base}}(t)}
                    {\sum_{t \in \mathcal{C}} \hat{y}_{\mathrm{base}}(t)^2},
  \label{eq:gain}
\end{equation}
and the calibrated base output is $\gamma \hat{y}_{\mathrm{base}}$.

The placement of the calibration segment requires care. Taking the first ten
seconds of each recording is not sufficient, because several recordings begin
with an extended rest period during which both the force and the base model
output are close to zero. Estimating \eqref{eq:gain} from such a segment is
dominated by noise and collapses the gain towards zero. We therefore locate the
segment by activity. The onset is the first sample at which the uncalibrated
base model output exceeds twenty percent of its 95th percentile for that
recording, and the calibration segment is the contiguous window of
\SI{10}{\second} beginning at that onset. The detector operates on the base model
output, which is a deterministic function of the sEMG signal, so the choice of
where to calibrate never consults the force channel. The window is additionally
constrained so that at least a quarter of the recording, and never fewer than one
full input window, remains available afterwards, which matters for the shorter
recordings. Estimated gains are clipped to $[10^{-3}, 10^{3}]$, and a gain that
reaches either bound is reported as a failed calibration rather than used
silently.

Two properties of this procedure should be stated explicitly. First, the force
channel is used inside the calibration segment, which is legitimate because a
reference contraction of known force is part of the deployment protocol rather
than information unavailable at test time. Second, every sample up to and
including the end of the calibration segment is excluded from evaluation. No
reported metric is computed on data that contributed to the calibration.

%------------------------------------------------------------------------------
\subsection{Residual Target}
\label{sec:target}

The quantity the network predicts is a scale-relative residual rather than the
raw error. For each recording we define
\begin{equation}
  \delta(t) \;=\; \frac{y(t) - \gamma \hat{y}_{\mathrm{base}}(t)}{\sigma},
  \qquad
  \sigma \;=\; \mathrm{p}_{95}\big\{\,\big|\gamma \hat{y}_{\mathrm{base}}\big|\,\big\},
  \label{eq:target}
\end{equation}
with $\sigma$ computed separately for each recording. At inference the network
output is multiplied by $\sigma$ before being added back in
\eqref{eq:decomposition}.

The reason for normalising per recording rather than fitting a single scaler on
the training set is specific to the cross-session setting. A global scaler stores
the mean and standard deviation of the residuals observed during training. Its
inverse transform then adds that stored mean to every prediction, including
predictions made on a new session whose residuals have a different magnitude and
frequently a different sign. The effect is a constant offset inherited from the
training sessions, applied to data the scaler was never fitted on. Defining the
target relative to $\sigma$ avoids this, because $\sigma$ is recomputed for each
new recording. Note also that $\sigma$ depends only on the base model output and
the calibration gain, both of which are available at deployment time.

%------------------------------------------------------------------------------
\subsection{Input Representation and Windowing}
\label{sec:inputs}

The network receives two input channels: the normalised sEMG $x_n$ and the
calibrated base model output divided by $\sigma$. Supplying the base model output
alongside the raw signal lets the network judge the correction in the context of
what the base model already predicted, rather than having to reconstruct that
prediction internally. Both channels are standardised by a scaler fitted on the
training recordings of the fold and applied unchanged to validation and test
recordings.

Inputs are formed as sliding windows of $L = 100$ samples, corresponding to
\SI{50}{\milli\second} at the \SI{2}{\kilo\hertz} sampling rate. Windows are
extracted with a stride of 20 samples, which reduces the strong redundancy
between neighbouring windows and keeps training time manageable without
discarding information that a stride of one would preserve. The network predicts
the residual at the final sample of each window, so no window contains
information from after the instant being predicted. Windows whose target sample
falls inside the calibration segment are discarded.

%------------------------------------------------------------------------------
\subsection{Compensation Architectures}
\label{sec:models}

Seven architectures are compared. Each maps a window
$\mathbf{X} \in \mathbb{R}^{L \times 2}$ to a scalar residual, so they are
interchangeable within the pipeline described above. Layer widths were chosen so
that parameter counts remain within one order of magnitude of each other
(Table~\ref{tab:models}). This matters for interpretation: a comparison in which
one model has far more capacity than the others measures capacity rather than
inductive bias. Every model is followed by the same prediction head, consisting
of layer normalisation, dropout, a linear layer to 16 units with a GELU
activation, and a final linear layer to a scalar.

\subsubsection{Convolutional BiLSTM Transformer}
\label{sec:cbt}

This is the reference architecture and combines three mechanisms. A
convolutional stem of two one-dimensional convolutions with kernel size three,
each followed by group normalisation and a GELU activation, extracts local
waveform morphology. A bidirectional LSTM with 48 hidden units per direction then
accumulates context in both temporal directions. The result is projected to a
48-dimensional representation, given sinusoidal positional encodings, and passed
through two pre-norm Transformer encoder layers with four attention heads and a
feedforward width of 96. Features are mean pooled over time before the head. The
model has 78{,}417 parameters. Its inclusion tests whether combining local,
recurrent and attention-based processing is beneficial for this task.

\subsubsection{Convolutional BiLSTM}
\label{sec:cbl}

This architecture removes the Transformer encoder and retains the convolutional
stem followed by a two-layer bidirectional LSTM with 64 hidden units per
direction. With 155{,}297 parameters it is the largest model in the comparison.
It isolates the contribution of self-attention: any difference relative to
Section~\ref{sec:cbt} is attributable to the attention stage rather than to the
convolutional or recurrent components.

\subsubsection{Long Short-Term Memory Network}
\label{sec:lstm}

A two-layer unidirectional LSTM with 96 hidden units, applied directly to the
input channels with no convolutional front end, provides a conventional
recurrent baseline. The hidden state at the final timestep is passed to the head.
The model has 114{,}657 parameters. Unlike the bidirectional variants it is
strictly causal within the window, which makes it the closest of the seven to a
model that could run in real time on a streaming signal.

\subsubsection{Transformer Encoder}
\label{sec:transformer}

This architecture uses attention alone. The input is projected linearly to 64
dimensions, given sinusoidal positional encodings, and processed by three
pre-norm encoder layers with four attention heads and a feedforward width of 128,
followed by mean pooling. It has 101{,}793 parameters. Because it has neither
recurrence nor convolution, all temporal structure must be recovered through
positional encoding and attention, which tests how much of the task depends on
explicitly sequential processing.

\subsubsection{Synergy State-Space Network}
\label{sec:synergy}

This architecture is motivated by the muscle synergy view of motor control, in
which descending commands are expressed as a small number of non-negative
activation primitives that are then filtered by muscle activation dynamics. The
input is projected onto $K = 16$ non-negative primitives and expanded to a bank
of first-order states,
\begin{align}
  \mathbf{u}(t) &= \mathrm{softplus}\!\big(\mathbf{W}\mathbf{x}(t)\big),
  \label{eq:synergy_u} \\
  \mathbf{h}(t) &= \mathbf{a} \odot \mathbf{h}(t-1)
                 + (\mathbf{1} - \mathbf{a}) \odot \mathbf{V}\mathbf{u}(t),
  \label{eq:synergy_h}
\end{align}
where $\mathbf{a} = \sigma(\boldsymbol{\theta}) \in (0,1)^{S}$ holds a learned
decay per state. The decays are initialised across a range of time constants so
that different states can specialise to fast and slow components of force
development. The state trajectory is mixed by a two-layer network before pooling.
The model has 70{,}481 parameters.

One limitation must be stated plainly. Classical synergy decomposition operates
on many simultaneously recorded EMG channels, and identifies groups of muscles
that are activated together. The present dataset provides a single sEMG channel,
so the layer in \eqref{eq:synergy_u} acts as a low-rank non-negative bottleneck
with associated activation dynamics rather than as a genuine multi-muscle synergy
extraction. We include it because the resulting inductive bias, namely a small
set of non-negative drives filtered by first-order dynamics, remains a reasonable
prior for this signal, but the architecture should not be read as a test of
synergy theory.

\subsubsection{Selective State-Space Model}
\label{sec:mamba}

This architecture follows the selective state-space formulation, in which the
discretisation step and the input and output projections of a diagonal state
space model are themselves functions of the input. The recurrence is
\begin{align}
  \mathbf{h}(t) &= \exp\!\big(\Delta_t \mathbf{A}\big)\,\mathbf{h}(t-1)
                 + \Delta_t \mathbf{B}_t\, x_t, \label{eq:mamba_h} \\
  y_t &= \mathbf{C}_t^{\!\top}\mathbf{h}(t) + D\, x_t, \label{eq:mamba_y}
\end{align}
where $\mathbf{A}$ is diagonal and negative, ensuring stability, and
$\Delta_t$, $\mathbf{B}_t$ and $\mathbf{C}_t$ are produced from the input at each
timestep. The block also contains a short depthwise causal convolution and a
gated residual branch. We use a model dimension of 64, a state dimension of 8,
and a single block, giving 30{,}770 parameters, the fewest of the seven models.
Input-dependent discretisation allows the effective time constant to vary with
the signal, which is attractive for a task in which the relevant timescale
changes between rapid onsets and slow sustained contractions.

We implement the selective scan directly in PyTorch rather than using the fused
CUDA kernels of the reference implementation. The computation is equivalent, but
the sequential scan over 100 timesteps is considerably slower, which is why a
single block is used. This affects training time only and not the reported
accuracy.

\subsubsection{Physics-Informed Neural Network}
\label{sec:pinn}

The final architecture adds an explicit dynamics prior to the loss. A
convolutional and bidirectional gated recurrent backbone emits a residual at
every timestep of the window rather than only at the last, which makes it
possible to evaluate a differential constraint across the window. Let $f(t)$
denote the reconstructed normalised force obtained by adding the predicted
residual sequence to the base model channel. The physics term penalises departure
from the same first-order relationship assumed by the base model in
\eqref{eq:lpf},
\begin{equation}
  \mathcal{L}_{\mathrm{phys}} \;=\;
  \Big\langle \Big( \tau \frac{\mathrm{d}f}{\mathrm{d}t} + f
                    - \big(w\, x_n + b\big) \Big)^{2} \Big\rangle,
  \label{eq:pinn}
\end{equation}
where the derivative is approximated by a first difference and $\tau$, $w$ and
$b$ are learned jointly with the network weights. The total objective is
$\mathcal{L}_{\mathrm{data}} + \lambda \mathcal{L}_{\mathrm{phys}}$. Only the
prediction at the final timestep contributes to the data term, so this model is
evaluated on exactly the same quantity as the other six. The model has 40{,}132
parameters. Because the inputs are standardised, the constraint holds only up to
an affine transformation of the input, which the learned $w$ and $b$ absorb.

\begin{table}[!t]
\caption{Compensation architectures compared. All share the same prediction
head and are trained with identical settings.}
\label{tab:models}
\centering
\begin{tabular}{llr}
\toprule
Model & Primary mechanism & Params \\
\midrule
Conv-BiLSTM-Transformer & convolution, recurrence, attention & 78{,}417 \\
Conv-BiLSTM             & convolution and recurrence         & 155{,}297 \\
LSTM                    & recurrence, causal                 & 114{,}657 \\
Transformer             & self-attention                     & 101{,}793 \\
Synergy-SSM             & non-negative bottleneck with SSM   & 70{,}481 \\
Selective SSM (Mamba)   & input-dependent state space        & 30{,}770 \\
PINN                    & dynamics-constrained residual      & 40{,}132 \\
\bottomrule
\end{tabular}
\end{table}

%------------------------------------------------------------------------------
\subsection{Training Protocol}
\label{sec:training}

All models are trained with identical settings. We use the AdamW optimiser with
a base learning rate of $3 \times 10^{-4}$ and weight decay of $10^{-4}$,
minimising a Huber loss on the normalised residual. The learning rate is warmed
up linearly over the first 500 optimiser steps and thereafter controlled by a
plateau scheduler that halves the rate after five epochs without improvement in
validation loss, down to a floor of $10^{-6}$. Gradients are clipped to a maximum
norm of one. Training runs for at most 100 epochs with early stopping after 15
epochs without improvement, and the weights from the epoch with the lowest
validation loss are restored before evaluation. The batch size is 64, and the
final incomplete batch of each epoch is dropped to avoid degenerate normalisation
statistics.

Several guards protect against silent failure. Batches that produce a non-finite
loss, or whose gradient norm is not finite, are skipped rather than applied, and
training aborts if more than twenty percent of the batches in an epoch are
skipped. Training also aborts if an epoch yields no finite validation loss, and
checkpoints containing non-finite weights are refused rather than written.

Two aspects of the protocol exist specifically to make the architecture
comparison fair. First, within a fold all seven models are trained on byte
identical inputs: the base model, the calibration gains, the feature scaler and
the extracted windows are all computed once and shared. Second, the random seed
is reset to the same value before each model is constructed, so differences
between models cannot be attributed to initialisation or to data ordering. The
only variable that changes between the seven runs is the architecture itself.

%------------------------------------------------------------------------------
\subsection{Evaluation Protocol}
\label{sec:protocol}

We evaluate under two protocols that answer different questions.

\subsubsection{Within-day evaluation}

Recordings from a single session are partitioned into training, validation and
test subsets, and the model never sees data from any other session. This measures
whether the residual can be learned at all when the recording conditions are held
fixed, and provides a reference point that is free of between-session variation.
Splits are performed at the level of whole recordings. This is essential rather
than conventional: consecutive samples in these signals are correlated at
approximately 0.99, so dividing a single recording into training and test
portions would place heavily overlapping windows on both sides of the boundary
and the reported error would be optimistic. Three rotating splits are run per
session so that different recordings serve as the test set, and results are
averaged across them.

\subsubsection{Cross-day evaluation}

One session is held out entirely for testing, a second is used for validation,
and the remaining sessions are used for training. The procedure is repeated so
that each eligible session serves once as the test session. This matches the
deployment scenario, in which a model trained on previously recorded sessions is
applied to a new electrode setup.

The validation session is selected by a fixed rotation among the eligible
sessions. We deliberately do not select it by similarity to the test session.
Doing so would require knowing the test session's force statistics, which are
precisely what is unavailable when a model must be chosen for deployment, and
would make the reported performance optimistic.

One session is designated as training data only and never serves as a test or
validation session. Its gain lies a factor of roughly two below that of every
other session, so holding it out would measure extrapolation beyond the range of
gains seen in training rather than transfer between comparable sessions. Its
value lies in training, where it widens the range of conditions the model has
observed. The extrapolation case is informative in its own right and is reported
separately in Section~\ref{sec:extrapolation}.

%------------------------------------------------------------------------------
\subsection{Metrics}
\label{sec:metrics}

We report root mean squared error (RMSE), mean absolute error (MAE) and the
coefficient of determination ($R^2$), each computed per recording and then
averaged. Every metric is computed on the evaluation segment only, that is, after
the calibration segment has been excluded. For every result we report the
corresponding base model performance on exactly the same samples, so that the
contribution of the learned compensation is always visible against the physics
based baseline rather than in isolation.

RMSE is treated as the primary metric. The reason concerns the behaviour of
$R^2$ on this dataset and should be understood before reading the results. The
coefficient of determination measures error relative to the variance of the
target within each recording. Several recordings contain long low-activity
intervals and consequently have small force variance, so a modest absolute error
produces a large negative $R^2$. Averaging such values across recordings yields
aggregate figures that are negative even where absolute error is clearly reduced,
and those figures are dominated by the least informative recordings rather than
by overall accuracy. We therefore rank methods by RMSE and MAE, which are
expressed in the units of the measured force and are comparable across
recordings, and report $R^2$ alongside them for completeness rather than as the
basis for any conclusion.

% Compile main.tex, not this file.
%==============================================================================

%==============================================================================
\section{Methodology}
\label{sec:method}

%------------------------------------------------------------------------------
\subsection{Problem Formulation}
\label{sec:formulation}

Let $x(t)$ denote the preprocessed surface electromyography (sEMG) envelope and
let $y(t)$ denote the simultaneously measured force. A physics-based base model
produces an estimate $\hat{y}_{\mathrm{base}}(t)$ from the sEMG signal alone.
This estimate captures the behaviour that the assumed physiology predicts, but
it leaves a systematic residual. We therefore write the final force estimate as
the sum of the base model output and a learned correction,
\begin{equation}
  \hat{y}(t) \;=\; \hat{y}_{\mathrm{base}}(t) \;+\; \Delta(t),
  \label{eq:decomposition}
\end{equation}
where $\Delta(t)$ is produced by a neural network from a window of recent
observations.

This decomposition is deliberate. The base model carries the part of the
relationship that is known in advance and that transfers between recording
sessions, while the network is asked only to explain what the base model cannot.
Compared with training a network to predict force directly, this keeps the
learned component small, makes its contribution measurable against an explicit
baseline, and ensures that a failure of the network degrades the estimate
towards the physics-based prediction rather than towards an arbitrary one.

%------------------------------------------------------------------------------
\subsection{Dataset}
\label{sec:dataset}

The dataset comprises 42 recordings collected across five sessions. Each
recording contains a time channel, a rectified sEMG channel, a low-pass filtered
sEMG channel, and a measured force channel, sampled at \SI{2}{\kilo\hertz}
($T_s = \SI{0.5}{\milli\second}$). This work uses the low-pass filtered sEMG
channel as the model input and the force channel as the target. Recordings in
four of the sessions are \SI{25}{\second} long (50{,}000 samples) and those in
the remaining session are \SI{15}{\second} long (30{,}000 samples).

Throughout this paper a session refers to a single electrode setup rather than
simply a calendar date. Within one session the electrodes were applied once and
left in place, so all recordings in that session share the same electrode
position, skin preparation and contact impedance. Between sessions the entire
setup was renewed. This distinction matters because the relationship between
sEMG amplitude and produced force is sensitive to exactly those factors.

Table~\ref{tab:sessions} summarises the sessions. For each recording we
characterise the input to output relationship by a single gain, defined as the
ratio of the 99th percentile of force to the 99th percentile of normalised sEMG.
The per-session median of this quantity varies from 0.44 to 5.86, a factor of
roughly thirteen between the extreme sessions, and the ordering is close to
monotonic across the five sessions. Variation of this magnitude means that a
model fitted on one session will systematically over- or under-predict on
another unless the gain is explicitly accounted for. Section~\ref{sec:calibration}
describes how we handle it.

\begin{table}[!t]
\caption{Recording sessions. The gain is the ratio of the 99th percentile force
to the 99th percentile normalised sEMG, reported as the per-session median.}
\label{tab:sessions}
\centering
\begin{tabular}{lrrrrl}
\toprule
Session & Files & Dur.~(s) & $F_{99}$ & Gain & Role \\
\midrule
1st\_Comp  & 5  & 25 & 3.68 & 3.70 & test / val / train \\
2nd\_Comp  & 9  & 25 & 5.56 & 5.86 & test / val / train \\
3rd\_Base  & 5  & 25 & 0.82 & 0.83 & test / val / train \\
4th\_Valid & 13 & 15 & 1.11 & 1.12 & test / val / train \\
5th\_Comp  & 10 & 25 & 0.44 & 0.44 & training only \\
\bottomrule
\end{tabular}
\end{table}

%------------------------------------------------------------------------------
\subsection{Per-Recording sEMG Normalisation}
\label{sec:semgnorm}

Each recording is normalised by its own amplitude statistics before entering the
model. We divide by the 99th percentile of the positive sEMG samples in that
recording,
\begin{equation}
  x_n(t) \;=\; \frac{x(t)}
                    {\mathrm{p}_{99}\{\,x(\tau) \;:\; x(\tau) > 0\,\}} ,
  \label{eq:semgnorm}
\end{equation}
using a percentile rather than the maximum so that isolated spikes do not set
the scale.

This step removes the component of between-session variation that originates in
electrode gain and contact quality. It is important to note that the
normalisation constant is computed from the input signal only. The sEMG signal
is available at deployment time, so the same quantity can be computed for a new
recording without any knowledge of the force that is being predicted. The
normalisation therefore introduces no label leakage, and this remains true when
the recording comes from a session the model has never seen.

%------------------------------------------------------------------------------
\subsection{Base Model}
\label{sec:basemodel}

The base model combines a static saturating nonlinearity with first-order
dynamics. The nonlinearity represents the compressive relationship between
muscle activation and produced force at higher activation levels, and the
dynamics represent the delay and smoothing introduced by the electromechanical
coupling between electrical activity and measurable force. The two stages are
\begin{align}
  z(t) &= g \,\tanh\!\big(s\, x_n(t)\big), \label{eq:tanh} \\
  G(\mathrm{s}) &= \frac{\omega_c}{\mathrm{s} + \omega_c},
                   \qquad \omega_c = 2\pi f_c, \label{eq:lpf}
\end{align}
where $g$ is a gain, $s$ shapes the saturation, and $f_c$ is the cutoff
frequency of the first-order stage. The continuous-time filter in
\eqref{eq:lpf} is discretised with the bilinear transform. Writing
$K = 2/T_s$, the resulting difference equation has numerator coefficients
$b_0 = b_1 = \omega_c/(K + \omega_c)$ and denominator coefficient
$a_1 = (\omega_c - K)/(K + \omega_c)$, and the base model output is
$\hat{y}_{\mathrm{base}} = G\{z\}$.

Including the first-order stage is necessary rather than cosmetic. When force is
plotted against sEMG for a single recording, the trajectory forms a hysteresis
loop: the force observed at a given sEMG level differs depending on whether the
contraction is building or releasing (Fig.~\ref{fig:hysteresis}). A static
nonlinearity is single-valued by construction and cannot reproduce this
behaviour, no matter how its parameters are chosen. The ablation in
Section~\ref{sec:ablation-lpf} quantifies what is lost when the stage is removed.

\begin{figure}[!t]
\centering
% \includegraphics[width=\columnwidth]{force_vs_semg_hysteresis.png}
\caption{Force plotted against normalised sEMG for representative recordings.
The trajectories form hysteresis loops, so the force observed at a given sEMG
level differs between the rising and falling phases of a contraction. A static
nonlinearity is single-valued and cannot represent this, which motivates the
first-order stage in \eqref{eq:lpf}.}
\label{fig:hysteresis}
\end{figure}

The three parameters $\{g, s, f_c\}$ are estimated by nonlinear least squares
using a trust region reflective algorithm, with $g$ and $s$ constrained to
$[0, 50]$ and $f_c$ constrained to $[1, 200]\,\si{\hertz}$. The fit is performed
on the concatenated residuals of all training recordings of a given fold. No
recording that will later be used for validation or testing contributes to the
fit, and the base model is refitted independently for each fold so that this
holds throughout.

%------------------------------------------------------------------------------
\subsection{Per-Session Gain Calibration}
\label{sec:calibration}

Because the input to output gain varies substantially across sessions, a single
fitted base model cannot serve all of them. We therefore estimate one scalar
gain per recording from a short calibration segment, which mirrors the practical
protocol of performing a brief reference contraction when electrodes are fitted.

Let $\mathcal{C}$ denote the set of samples in the calibration segment. The
calibration gain is the least squares scaling that best aligns the base model
output with the measured force over that segment,
\begin{equation}
  \gamma \;=\; \frac{\sum_{t \in \mathcal{C}} y(t)\, \hat{y}_{\mathrm{base}}(t)}
                    {\sum_{t \in \mathcal{C}} \hat{y}_{\mathrm{base}}(t)^2},
  \label{eq:gain}
\end{equation}
and the calibrated base output is $\gamma \hat{y}_{\mathrm{base}}$.

The placement of the calibration segment requires care. Taking the first ten
seconds of each recording is not sufficient, because several recordings begin
with an extended rest period during which both the force and the base model
output are close to zero. Estimating \eqref{eq:gain} from such a segment is
dominated by noise and collapses the gain towards zero. We therefore locate the
segment by activity. The onset is the first sample at which the uncalibrated
base model output exceeds twenty percent of its 95th percentile for that
recording, and the calibration segment is the contiguous window of
\SI{10}{\second} beginning at that onset. The detector operates on the base model
output, which is a deterministic function of the sEMG signal, so the choice of
where to calibrate never consults the force channel. The window is additionally
constrained so that at least a quarter of the recording, and never fewer than one
full input window, remains available afterwards, which matters for the shorter
recordings. Estimated gains are clipped to $[10^{-3}, 10^{3}]$, and a gain that
reaches either bound is reported as a failed calibration rather than used
silently.

Two properties of this procedure should be stated explicitly. First, the force
channel is used inside the calibration segment, which is legitimate because a
reference contraction of known force is part of the deployment protocol rather
than information unavailable at test time. Second, every sample up to and
including the end of the calibration segment is excluded from evaluation. No
reported metric is computed on data that contributed to the calibration.

%------------------------------------------------------------------------------
\subsection{Residual Target}
\label{sec:target}

The quantity the network predicts is a scale-relative residual rather than the
raw error. For each recording we define
\begin{equation}
  \delta(t) \;=\; \frac{y(t) - \gamma \hat{y}_{\mathrm{base}}(t)}{\sigma},
  \qquad
  \sigma \;=\; \mathrm{p}_{95}\big\{\,\big|\gamma \hat{y}_{\mathrm{base}}\big|\,\big\},
  \label{eq:target}
\end{equation}
with $\sigma$ computed separately for each recording. At inference the network
output is multiplied by $\sigma$ before being added back in
\eqref{eq:decomposition}.

The reason for normalising per recording rather than fitting a single scaler on
the training set is specific to the cross-session setting. A global scaler stores
the mean and standard deviation of the residuals observed during training. Its
inverse transform then adds that stored mean to every prediction, including
predictions made on a new session whose residuals have a different magnitude and
frequently a different sign. The effect is a constant offset inherited from the
training sessions, applied to data the scaler was never fitted on. Defining the
target relative to $\sigma$ avoids this, because $\sigma$ is recomputed for each
new recording. Note also that $\sigma$ depends only on the base model output and
the calibration gain, both of which are available at deployment time.

%------------------------------------------------------------------------------
\subsection{Input Representation and Windowing}
\label{sec:inputs}

The network receives two input channels: the normalised sEMG $x_n$ and the
calibrated base model output divided by $\sigma$. Supplying the base model output
alongside the raw signal lets the network judge the correction in the context of
what the base model already predicted, rather than having to reconstruct that
prediction internally. Both channels are standardised by a scaler fitted on the
training recordings of the fold and applied unchanged to validation and test
recordings.

Inputs are formed as sliding windows of $L = 100$ samples, corresponding to
\SI{50}{\milli\second} at the \SI{2}{\kilo\hertz} sampling rate. Windows are
extracted with a stride of 20 samples, which reduces the strong redundancy
between neighbouring windows and keeps training time manageable without
discarding information that a stride of one would preserve. The network predicts
the residual at the final sample of each window, so no window contains
information from after the instant being predicted. Windows whose target sample
falls inside the calibration segment are discarded.

%------------------------------------------------------------------------------
\subsection{Compensation Architectures}
\label{sec:models}

Seven architectures are compared. Each maps a window
$\mathbf{X} \in \mathbb{R}^{L \times 2}$ to a scalar residual, so they are
interchangeable within the pipeline described above. Layer widths were chosen so
that parameter counts remain within one order of magnitude of each other
(Table~\ref{tab:models}). This matters for interpretation: a comparison in which
one model has far more capacity than the others measures capacity rather than
inductive bias. Every model is followed by the same prediction head, consisting
of layer normalisation, dropout, a linear layer to 16 units with a GELU
activation, and a final linear layer to a scalar.

\subsubsection{Convolutional BiLSTM Transformer}
\label{sec:cbt}

This is the reference architecture and combines three mechanisms. A
convolutional stem of two one-dimensional convolutions with kernel size three,
each followed by group normalisation and a GELU activation, extracts local
waveform morphology. A bidirectional LSTM with 48 hidden units per direction then
accumulates context in both temporal directions. The result is projected to a
48-dimensional representation, given sinusoidal positional encodings, and passed
through two pre-norm Transformer encoder layers with four attention heads and a
feedforward width of 96. Features are mean pooled over time before the head. The
model has 78{,}417 parameters. Its inclusion tests whether combining local,
recurrent and attention-based processing is beneficial for this task.

\subsubsection{Convolutional BiLSTM}
\label{sec:cbl}

This architecture removes the Transformer encoder and retains the convolutional
stem followed by a two-layer bidirectional LSTM with 64 hidden units per
direction. With 155{,}297 parameters it is the largest model in the comparison.
It isolates the contribution of self-attention: any difference relative to
Section~\ref{sec:cbt} is attributable to the attention stage rather than to the
convolutional or recurrent components.

\subsubsection{Long Short-Term Memory Network}
\label{sec:lstm}

A two-layer unidirectional LSTM with 96 hidden units, applied directly to the
input channels with no convolutional front end, provides a conventional
recurrent baseline. The hidden state at the final timestep is passed to the head.
The model has 114{,}657 parameters. Unlike the bidirectional variants it is
strictly causal within the window, which makes it the closest of the seven to a
model that could run in real time on a streaming signal.

\subsubsection{Transformer Encoder}
\label{sec:transformer}

This architecture uses attention alone. The input is projected linearly to 64
dimensions, given sinusoidal positional encodings, and processed by three
pre-norm encoder layers with four attention heads and a feedforward width of 128,
followed by mean pooling. It has 101{,}793 parameters. Because it has neither
recurrence nor convolution, all temporal structure must be recovered through
positional encoding and attention, which tests how much of the task depends on
explicitly sequential processing.

\subsubsection{Synergy State-Space Network}
\label{sec:synergy}

This architecture is motivated by the muscle synergy view of motor control, in
which descending commands are expressed as a small number of non-negative
activation primitives that are then filtered by muscle activation dynamics. The
input is projected onto $K = 16$ non-negative primitives and expanded to a bank
of first-order states,
\begin{align}
  \mathbf{u}(t) &= \mathrm{softplus}\!\big(\mathbf{W}\mathbf{x}(t)\big),
  \label{eq:synergy_u} \\
  \mathbf{h}(t) &= \mathbf{a} \odot \mathbf{h}(t-1)
                 + (\mathbf{1} - \mathbf{a}) \odot \mathbf{V}\mathbf{u}(t),
  \label{eq:synergy_h}
\end{align}
where $\mathbf{a} = \sigma(\boldsymbol{\theta}) \in (0,1)^{S}$ holds a learned
decay per state. The decays are initialised across a range of time constants so
that different states can specialise to fast and slow components of force
development. The state trajectory is mixed by a two-layer network before pooling.
The model has 70{,}481 parameters.

One limitation must be stated plainly. Classical synergy decomposition operates
on many simultaneously recorded EMG channels, and identifies groups of muscles
that are activated together. The present dataset provides a single sEMG channel,
so the layer in \eqref{eq:synergy_u} acts as a low-rank non-negative bottleneck
with associated activation dynamics rather than as a genuine multi-muscle synergy
extraction. We include it because the resulting inductive bias, namely a small
set of non-negative drives filtered by first-order dynamics, remains a reasonable
prior for this signal, but the architecture should not be read as a test of
synergy theory.

\subsubsection{Selective State-Space Model}
\label{sec:mamba}

This architecture follows the selective state-space formulation, in which the
discretisation step and the input and output projections of a diagonal state
space model are themselves functions of the input. The recurrence is
\begin{align}
  \mathbf{h}(t) &= \exp\!\big(\Delta_t \mathbf{A}\big)\,\mathbf{h}(t-1)
                 + \Delta_t \mathbf{B}_t\, x_t, \label{eq:mamba_h} \\
  y_t &= \mathbf{C}_t^{\!\top}\mathbf{h}(t) + D\, x_t, \label{eq:mamba_y}
\end{align}
where $\mathbf{A}$ is diagonal and negative, ensuring stability, and
$\Delta_t$, $\mathbf{B}_t$ and $\mathbf{C}_t$ are produced from the input at each
timestep. The block also contains a short depthwise causal convolution and a
gated residual branch. We use a model dimension of 64, a state dimension of 8,
and a single block, giving 30{,}770 parameters, the fewest of the seven models.
Input-dependent discretisation allows the effective time constant to vary with
the signal, which is attractive for a task in which the relevant timescale
changes between rapid onsets and slow sustained contractions.

We implement the selective scan directly in PyTorch rather than using the fused
CUDA kernels of the reference implementation. The computation is equivalent, but
the sequential scan over 100 timesteps is considerably slower, which is why a
single block is used. This affects training time only and not the reported
accuracy.

\subsubsection{Physics-Informed Neural Network}
\label{sec:pinn}

The final architecture adds an explicit dynamics prior to the loss. A
convolutional and bidirectional gated recurrent backbone emits a residual at
every timestep of the window rather than only at the last, which makes it
possible to evaluate a differential constraint across the window. Let $f(t)$
denote the reconstructed normalised force obtained by adding the predicted
residual sequence to the base model channel. The physics term penalises departure
from the same first-order relationship assumed by the base model in
\eqref{eq:lpf},
\begin{equation}
  \mathcal{L}_{\mathrm{phys}} \;=\;
  \Big\langle \Big( \tau \frac{\mathrm{d}f}{\mathrm{d}t} + f
                    - \big(w\, x_n + b\big) \Big)^{2} \Big\rangle,
  \label{eq:pinn}
\end{equation}
where the derivative is approximated by a first difference and $\tau$, $w$ and
$b$ are learned jointly with the network weights. The total objective is
$\mathcal{L}_{\mathrm{data}} + \lambda \mathcal{L}_{\mathrm{phys}}$. Only the
prediction at the final timestep contributes to the data term, so this model is
evaluated on exactly the same quantity as the other six. The model has 40{,}132
parameters. Because the inputs are standardised, the constraint holds only up to
an affine transformation of the input, which the learned $w$ and $b$ absorb.

\begin{table}[!t]
\caption{Compensation architectures compared. All share the same prediction
head and are trained with identical settings.}
\label{tab:models}
\centering
\begin{tabular}{llr}
\toprule
Model & Primary mechanism & Params \\
\midrule
Conv-BiLSTM-Transformer & convolution, recurrence, attention & 78{,}417 \\
Conv-BiLSTM             & convolution and recurrence         & 155{,}297 \\
LSTM                    & recurrence, causal                 & 114{,}657 \\
Transformer             & self-attention                     & 101{,}793 \\
Synergy-SSM             & non-negative bottleneck with SSM   & 70{,}481 \\
Selective SSM (Mamba)   & input-dependent state space        & 30{,}770 \\
PINN                    & dynamics-constrained residual      & 40{,}132 \\
\bottomrule
\end{tabular}
\end{table}

%------------------------------------------------------------------------------
\subsection{Training Protocol}
\label{sec:training}

All models are trained with identical settings. We use the AdamW optimiser with
a base learning rate of $3 \times 10^{-4}$ and weight decay of $10^{-4}$,
minimising a Huber loss on the normalised residual. The learning rate is warmed
up linearly over the first 500 optimiser steps and thereafter controlled by a
plateau scheduler that halves the rate after five epochs without improvement in
validation loss, down to a floor of $10^{-6}$. Gradients are clipped to a maximum
norm of one. Training runs for at most 100 epochs with early stopping after 15
epochs without improvement, and the weights from the epoch with the lowest
validation loss are restored before evaluation. The batch size is 64, and the
final incomplete batch of each epoch is dropped to avoid degenerate normalisation
statistics.

Several guards protect against silent failure. Batches that produce a non-finite
loss, or whose gradient norm is not finite, are skipped rather than applied, and
training aborts if more than twenty percent of the batches in an epoch are
skipped. Training also aborts if an epoch yields no finite validation loss, and
checkpoints containing non-finite weights are refused rather than written.

Two aspects of the protocol exist specifically to make the architecture
comparison fair. First, within a fold all seven models are trained on byte
identical inputs: the base model, the calibration gains, the feature scaler and
the extracted windows are all computed once and shared. Second, the random seed
is reset to the same value before each model is constructed, so differences
between models cannot be attributed to initialisation or to data ordering. The
only variable that changes between the seven runs is the architecture itself.

%------------------------------------------------------------------------------
\subsection{Evaluation Protocol}
\label{sec:protocol}

We evaluate under two protocols that answer different questions.

\subsubsection{Within-day evaluation}

Recordings from a single session are partitioned into training, validation and
test subsets, and the model never sees data from any other session. This measures
whether the residual can be learned at all when the recording conditions are held
fixed, and provides a reference point that is free of between-session variation.
Splits are performed at the level of whole recordings. This is essential rather
than conventional: consecutive samples in these signals are correlated at
approximately 0.99, so dividing a single recording into training and test
portions would place heavily overlapping windows on both sides of the boundary
and the reported error would be optimistic. Three rotating splits are run per
session so that different recordings serve as the test set, and results are
averaged across them.

\subsubsection{Cross-day evaluation}

One session is held out entirely for testing, a second is used for validation,
and the remaining sessions are used for training. The procedure is repeated so
that each eligible session serves once as the test session. This matches the
deployment scenario, in which a model trained on previously recorded sessions is
applied to a new electrode setup.

The validation session is selected by a fixed rotation among the eligible
sessions. We deliberately do not select it by similarity to the test session.
Doing so would require knowing the test session's force statistics, which are
precisely what is unavailable when a model must be chosen for deployment, and
would make the reported performance optimistic.

One session is designated as training data only and never serves as a test or
validation session. Its gain lies a factor of roughly two below that of every
other session, so holding it out would measure extrapolation beyond the range of
gains seen in training rather than transfer between comparable sessions. Its
value lies in training, where it widens the range of conditions the model has
observed. The extrapolation case is informative in its own right and is reported
separately in Section~\ref{sec:extrapolation}.

%------------------------------------------------------------------------------
\subsection{Metrics}
\label{sec:metrics}

We report root mean squared error (RMSE), mean absolute error (MAE) and the
coefficient of determination ($R^2$), each computed per recording and then
averaged. Every metric is computed on the evaluation segment only, that is, after
the calibration segment has been excluded. For every result we report the
corresponding base model performance on exactly the same samples, so that the
contribution of the learned compensation is always visible against the physics
based baseline rather than in isolation.

RMSE is treated as the primary metric. The reason concerns the behaviour of
$R^2$ on this dataset and should be understood before reading the results. The
coefficient of determination measures error relative to the variance of the
target within each recording. Several recordings contain long low-activity
intervals and consequently have small force variance, so a modest absolute error
produces a large negative $R^2$. Averaging such values across recordings yields
aggregate figures that are negative even where absolute error is clearly reduced,
and those figures are dominated by the least informative recordings rather than
by overall accuracy. We therefore rank methods by RMSE and MAE, which are
expressed in the units of the measured force and are comparable across
recordings, and report $R^2$ alongside them for completeness rather than as the
basis for any conclusion.

% Compile main.tex, not this file.
%==============================================================================

%==============================================================================
\section{Methodology}
\label{sec:method}

%------------------------------------------------------------------------------
\subsection{Problem Formulation}
\label{sec:formulation}

Let $x(t)$ denote the preprocessed surface electromyography (sEMG) envelope and
let $y(t)$ denote the simultaneously measured force. A physics-based base model
produces an estimate $\hat{y}_{\mathrm{base}}(t)$ from the sEMG signal alone.
This estimate captures the behaviour that the assumed physiology predicts, but
it leaves a systematic residual. We therefore write the final force estimate as
the sum of the base model output and a learned correction,
\begin{equation}
  \hat{y}(t) \;=\; \hat{y}_{\mathrm{base}}(t) \;+\; \Delta(t),
  \label{eq:decomposition}
\end{equation}
where $\Delta(t)$ is produced by a neural network from a window of recent
observations.

This decomposition is deliberate. The base model carries the part of the
relationship that is known in advance and that transfers between recording
sessions, while the network is asked only to explain what the base model cannot.
Compared with training a network to predict force directly, this keeps the
learned component small, makes its contribution measurable against an explicit
baseline, and ensures that a failure of the network degrades the estimate
towards the physics-based prediction rather than towards an arbitrary one.

%------------------------------------------------------------------------------
\subsection{Dataset}
\label{sec:dataset}

The dataset comprises 42 recordings collected across five sessions. Each
recording contains a time channel, a rectified sEMG channel, a low-pass filtered
sEMG channel, and a measured force channel, sampled at \SI{2}{\kilo\hertz}
($T_s = \SI{0.5}{\milli\second}$). This work uses the low-pass filtered sEMG
channel as the model input and the force channel as the target. Recordings in
four of the sessions are \SI{25}{\second} long (50{,}000 samples) and those in
the remaining session are \SI{15}{\second} long (30{,}000 samples).

Throughout this paper a session refers to a single electrode setup rather than
simply a calendar date. Within one session the electrodes were applied once and
left in place, so all recordings in that session share the same electrode
position, skin preparation and contact impedance. Between sessions the entire
setup was renewed. This distinction matters because the relationship between
sEMG amplitude and produced force is sensitive to exactly those factors.

Table~\ref{tab:sessions} summarises the sessions. For each recording we
characterise the input to output relationship by a single gain, defined as the
ratio of the 99th percentile of force to the 99th percentile of normalised sEMG.
The per-session median of this quantity varies from 0.44 to 5.86, a factor of
roughly thirteen between the extreme sessions, and the ordering is close to
monotonic across the five sessions. Variation of this magnitude means that a
model fitted on one session will systematically over- or under-predict on
another unless the gain is explicitly accounted for. Section~\ref{sec:calibration}
describes how we handle it.

\begin{table}[!t]
\caption{Recording sessions. The gain is the ratio of the 99th percentile force
to the 99th percentile normalised sEMG, reported as the per-session median.}
\label{tab:sessions}
\centering
\begin{tabular}{lrrrrl}
\toprule
Session & Files & Dur.~(s) & $F_{99}$ & Gain & Role \\
\midrule
1st\_Comp  & 5  & 25 & 3.68 & 3.70 & test / val / train \\
2nd\_Comp  & 9  & 25 & 5.56 & 5.86 & test / val / train \\
3rd\_Base  & 5  & 25 & 0.82 & 0.83 & test / val / train \\
4th\_Valid & 13 & 15 & 1.11 & 1.12 & test / val / train \\
5th\_Comp  & 10 & 25 & 0.44 & 0.44 & training only \\
\bottomrule
\end{tabular}
\end{table}

%------------------------------------------------------------------------------
\subsection{Per-Recording sEMG Normalisation}
\label{sec:semgnorm}

Each recording is normalised by its own amplitude statistics before entering the
model. We divide by the 99th percentile of the positive sEMG samples in that
recording,
\begin{equation}
  x_n(t) \;=\; \frac{x(t)}
                    {\mathrm{p}_{99}\{\,x(\tau) \;:\; x(\tau) > 0\,\}} ,
  \label{eq:semgnorm}
\end{equation}
using a percentile rather than the maximum so that isolated spikes do not set
the scale.

This step removes the component of between-session variation that originates in
electrode gain and contact quality. It is important to note that the
normalisation constant is computed from the input signal only. The sEMG signal
is available at deployment time, so the same quantity can be computed for a new
recording without any knowledge of the force that is being predicted. The
normalisation therefore introduces no label leakage, and this remains true when
the recording comes from a session the model has never seen.

%------------------------------------------------------------------------------
\subsection{Base Model}
\label{sec:basemodel}

The base model combines a static saturating nonlinearity with first-order
dynamics. The nonlinearity represents the compressive relationship between
muscle activation and produced force at higher activation levels, and the
dynamics represent the delay and smoothing introduced by the electromechanical
coupling between electrical activity and measurable force. The two stages are
\begin{align}
  z(t) &= g \,\tanh\!\big(s\, x_n(t)\big), \label{eq:tanh} \\
  G(\mathrm{s}) &= \frac{\omega_c}{\mathrm{s} + \omega_c},
                   \qquad \omega_c = 2\pi f_c, \label{eq:lpf}
\end{align}
where $g$ is a gain, $s$ shapes the saturation, and $f_c$ is the cutoff
frequency of the first-order stage. The continuous-time filter in
\eqref{eq:lpf} is discretised with the bilinear transform. Writing
$K = 2/T_s$, the resulting difference equation has numerator coefficients
$b_0 = b_1 = \omega_c/(K + \omega_c)$ and denominator coefficient
$a_1 = (\omega_c - K)/(K + \omega_c)$, and the base model output is
$\hat{y}_{\mathrm{base}} = G\{z\}$.

Including the first-order stage is necessary rather than cosmetic. When force is
plotted against sEMG for a single recording, the trajectory forms a hysteresis
loop: the force observed at a given sEMG level differs depending on whether the
contraction is building or releasing (Fig.~\ref{fig:hysteresis}). A static
nonlinearity is single-valued by construction and cannot reproduce this
behaviour, no matter how its parameters are chosen. The ablation in
Section~\ref{sec:ablation-lpf} quantifies what is lost when the stage is removed.

\begin{figure}[!t]
\centering
% \includegraphics[width=\columnwidth]{force_vs_semg_hysteresis.png}
\caption{Force plotted against normalised sEMG for representative recordings.
The trajectories form hysteresis loops, so the force observed at a given sEMG
level differs between the rising and falling phases of a contraction. A static
nonlinearity is single-valued and cannot represent this, which motivates the
first-order stage in \eqref{eq:lpf}.}
\label{fig:hysteresis}
\end{figure}

The three parameters $\{g, s, f_c\}$ are estimated by nonlinear least squares
using a trust region reflective algorithm, with $g$ and $s$ constrained to
$[0, 50]$ and $f_c$ constrained to $[1, 200]\,\si{\hertz}$. The fit is performed
on the concatenated residuals of all training recordings of a given fold. No
recording that will later be used for validation or testing contributes to the
fit, and the base model is refitted independently for each fold so that this
holds throughout.

%------------------------------------------------------------------------------
\subsection{Per-Session Gain Calibration}
\label{sec:calibration}

Because the input to output gain varies substantially across sessions, a single
fitted base model cannot serve all of them. We therefore estimate one scalar
gain per recording from a short calibration segment, which mirrors the practical
protocol of performing a brief reference contraction when electrodes are fitted.

Let $\mathcal{C}$ denote the set of samples in the calibration segment. The
calibration gain is the least squares scaling that best aligns the base model
output with the measured force over that segment,
\begin{equation}
  \gamma \;=\; \frac{\sum_{t \in \mathcal{C}} y(t)\, \hat{y}_{\mathrm{base}}(t)}
                    {\sum_{t \in \mathcal{C}} \hat{y}_{\mathrm{base}}(t)^2},
  \label{eq:gain}
\end{equation}
and the calibrated base output is $\gamma \hat{y}_{\mathrm{base}}$.

The placement of the calibration segment requires care. Taking the first ten
seconds of each recording is not sufficient, because several recordings begin
with an extended rest period during which both the force and the base model
output are close to zero. Estimating \eqref{eq:gain} from such a segment is
dominated by noise and collapses the gain towards zero. We therefore locate the
segment by activity. The onset is the first sample at which the uncalibrated
base model output exceeds twenty percent of its 95th percentile for that
recording, and the calibration segment is the contiguous window of
\SI{10}{\second} beginning at that onset. The detector operates on the base model
output, which is a deterministic function of the sEMG signal, so the choice of
where to calibrate never consults the force channel. The window is additionally
constrained so that at least a quarter of the recording, and never fewer than one
full input window, remains available afterwards, which matters for the shorter
recordings. Estimated gains are clipped to $[10^{-3}, 10^{3}]$, and a gain that
reaches either bound is reported as a failed calibration rather than used
silently.

Two properties of this procedure should be stated explicitly. First, the force
channel is used inside the calibration segment, which is legitimate because a
reference contraction of known force is part of the deployment protocol rather
than information unavailable at test time. Second, every sample up to and
including the end of the calibration segment is excluded from evaluation. No
reported metric is computed on data that contributed to the calibration.

%------------------------------------------------------------------------------
\subsection{Residual Target}
\label{sec:target}

The quantity the network predicts is a scale-relative residual rather than the
raw error. For each recording we define
\begin{equation}
  \delta(t) \;=\; \frac{y(t) - \gamma \hat{y}_{\mathrm{base}}(t)}{\sigma},
  \qquad
  \sigma \;=\; \mathrm{p}_{95}\big\{\,\big|\gamma \hat{y}_{\mathrm{base}}\big|\,\big\},
  \label{eq:target}
\end{equation}
with $\sigma$ computed separately for each recording. At inference the network
output is multiplied by $\sigma$ before being added back in
\eqref{eq:decomposition}.

The reason for normalising per recording rather than fitting a single scaler on
the training set is specific to the cross-session setting. A global scaler stores
the mean and standard deviation of the residuals observed during training. Its
inverse transform then adds that stored mean to every prediction, including
predictions made on a new session whose residuals have a different magnitude and
frequently a different sign. The effect is a constant offset inherited from the
training sessions, applied to data the scaler was never fitted on. Defining the
target relative to $\sigma$ avoids this, because $\sigma$ is recomputed for each
new recording. Note also that $\sigma$ depends only on the base model output and
the calibration gain, both of which are available at deployment time.

%------------------------------------------------------------------------------
\subsection{Input Representation and Windowing}
\label{sec:inputs}

The network receives two input channels: the normalised sEMG $x_n$ and the
calibrated base model output divided by $\sigma$. Supplying the base model output
alongside the raw signal lets the network judge the correction in the context of
what the base model already predicted, rather than having to reconstruct that
prediction internally. Both channels are standardised by a scaler fitted on the
training recordings of the fold and applied unchanged to validation and test
recordings.

Inputs are formed as sliding windows of $L = 100$ samples, corresponding to
\SI{50}{\milli\second} at the \SI{2}{\kilo\hertz} sampling rate. Windows are
extracted with a stride of 20 samples, which reduces the strong redundancy
between neighbouring windows and keeps training time manageable without
discarding information that a stride of one would preserve. The network predicts
the residual at the final sample of each window, so no window contains
information from after the instant being predicted. Windows whose target sample
falls inside the calibration segment are discarded.

%------------------------------------------------------------------------------
\subsection{Compensation Architectures}
\label{sec:models}

Seven architectures are compared. Each maps a window
$\mathbf{X} \in \mathbb{R}^{L \times 2}$ to a scalar residual, so they are
interchangeable within the pipeline described above. Layer widths were chosen so
that parameter counts remain within one order of magnitude of each other
(Table~\ref{tab:models}). This matters for interpretation: a comparison in which
one model has far more capacity than the others measures capacity rather than
inductive bias. Every model is followed by the same prediction head, consisting
of layer normalisation, dropout, a linear layer to 16 units with a GELU
activation, and a final linear layer to a scalar.

\subsubsection{Convolutional BiLSTM Transformer}
\label{sec:cbt}

This is the reference architecture and combines three mechanisms. A
convolutional stem of two one-dimensional convolutions with kernel size three,
each followed by group normalisation and a GELU activation, extracts local
waveform morphology. A bidirectional LSTM with 48 hidden units per direction then
accumulates context in both temporal directions. The result is projected to a
48-dimensional representation, given sinusoidal positional encodings, and passed
through two pre-norm Transformer encoder layers with four attention heads and a
feedforward width of 96. Features are mean pooled over time before the head. The
model has 78{,}417 parameters. Its inclusion tests whether combining local,
recurrent and attention-based processing is beneficial for this task.

\subsubsection{Convolutional BiLSTM}
\label{sec:cbl}

This architecture removes the Transformer encoder and retains the convolutional
stem followed by a two-layer bidirectional LSTM with 64 hidden units per
direction. With 155{,}297 parameters it is the largest model in the comparison.
It isolates the contribution of self-attention: any difference relative to
Section~\ref{sec:cbt} is attributable to the attention stage rather than to the
convolutional or recurrent components.

\subsubsection{Long Short-Term Memory Network}
\label{sec:lstm}

A two-layer unidirectional LSTM with 96 hidden units, applied directly to the
input channels with no convolutional front end, provides a conventional
recurrent baseline. The hidden state at the final timestep is passed to the head.
The model has 114{,}657 parameters. Unlike the bidirectional variants it is
strictly causal within the window, which makes it the closest of the seven to a
model that could run in real time on a streaming signal.

\subsubsection{Transformer Encoder}
\label{sec:transformer}

This architecture uses attention alone. The input is projected linearly to 64
dimensions, given sinusoidal positional encodings, and processed by three
pre-norm encoder layers with four attention heads and a feedforward width of 128,
followed by mean pooling. It has 101{,}793 parameters. Because it has neither
recurrence nor convolution, all temporal structure must be recovered through
positional encoding and attention, which tests how much of the task depends on
explicitly sequential processing.

\subsubsection{Synergy State-Space Network}
\label{sec:synergy}

This architecture is motivated by the muscle synergy view of motor control, in
which descending commands are expressed as a small number of non-negative
activation primitives that are then filtered by muscle activation dynamics. The
input is projected onto $K = 16$ non-negative primitives and expanded to a bank
of first-order states,
\begin{align}
  \mathbf{u}(t) &= \mathrm{softplus}\!\big(\mathbf{W}\mathbf{x}(t)\big),
  \label{eq:synergy_u} \\
  \mathbf{h}(t) &= \mathbf{a} \odot \mathbf{h}(t-1)
                 + (\mathbf{1} - \mathbf{a}) \odot \mathbf{V}\mathbf{u}(t),
  \label{eq:synergy_h}
\end{align}
where $\mathbf{a} = \sigma(\boldsymbol{\theta}) \in (0,1)^{S}$ holds a learned
decay per state. The decays are initialised across a range of time constants so
that different states can specialise to fast and slow components of force
development. The state trajectory is mixed by a two-layer network before pooling.
The model has 70{,}481 parameters.

One limitation must be stated plainly. Classical synergy decomposition operates
on many simultaneously recorded EMG channels, and identifies groups of muscles
that are activated together. The present dataset provides a single sEMG channel,
so the layer in \eqref{eq:synergy_u} acts as a low-rank non-negative bottleneck
with associated activation dynamics rather than as a genuine multi-muscle synergy
extraction. We include it because the resulting inductive bias, namely a small
set of non-negative drives filtered by first-order dynamics, remains a reasonable
prior for this signal, but the architecture should not be read as a test of
synergy theory.

\subsubsection{Selective State-Space Model}
\label{sec:mamba}

This architecture follows the selective state-space formulation, in which the
discretisation step and the input and output projections of a diagonal state
space model are themselves functions of the input. The recurrence is
\begin{align}
  \mathbf{h}(t) &= \exp\!\big(\Delta_t \mathbf{A}\big)\,\mathbf{h}(t-1)
                 + \Delta_t \mathbf{B}_t\, x_t, \label{eq:mamba_h} \\
  y_t &= \mathbf{C}_t^{\!\top}\mathbf{h}(t) + D\, x_t, \label{eq:mamba_y}
\end{align}
where $\mathbf{A}$ is diagonal and negative, ensuring stability, and
$\Delta_t$, $\mathbf{B}_t$ and $\mathbf{C}_t$ are produced from the input at each
timestep. The block also contains a short depthwise causal convolution and a
gated residual branch. We use a model dimension of 64, a state dimension of 8,
and a single block, giving 30{,}770 parameters, the fewest of the seven models.
Input-dependent discretisation allows the effective time constant to vary with
the signal, which is attractive for a task in which the relevant timescale
changes between rapid onsets and slow sustained contractions.

We implement the selective scan directly in PyTorch rather than using the fused
CUDA kernels of the reference implementation. The computation is equivalent, but
the sequential scan over 100 timesteps is considerably slower, which is why a
single block is used. This affects training time only and not the reported
accuracy.

\subsubsection{Physics-Informed Neural Network}
\label{sec:pinn}

The final architecture adds an explicit dynamics prior to the loss. A
convolutional and bidirectional gated recurrent backbone emits a residual at
every timestep of the window rather than only at the last, which makes it
possible to evaluate a differential constraint across the window. Let $f(t)$
denote the reconstructed normalised force obtained by adding the predicted
residual sequence to the base model channel. The physics term penalises departure
from the same first-order relationship assumed by the base model in
\eqref{eq:lpf},
\begin{equation}
  \mathcal{L}_{\mathrm{phys}} \;=\;
  \Big\langle \Big( \tau \frac{\mathrm{d}f}{\mathrm{d}t} + f
                    - \big(w\, x_n + b\big) \Big)^{2} \Big\rangle,
  \label{eq:pinn}
\end{equation}
where the derivative is approximated by a first difference and $\tau$, $w$ and
$b$ are learned jointly with the network weights. The total objective is
$\mathcal{L}_{\mathrm{data}} + \lambda \mathcal{L}_{\mathrm{phys}}$. Only the
prediction at the final timestep contributes to the data term, so this model is
evaluated on exactly the same quantity as the other six. The model has 40{,}132
parameters. Because the inputs are standardised, the constraint holds only up to
an affine transformation of the input, which the learned $w$ and $b$ absorb.

\begin{table}[!t]
\caption{Compensation architectures compared. All share the same prediction
head and are trained with identical settings.}
\label{tab:models}
\centering
\begin{tabular}{llr}
\toprule
Model & Primary mechanism & Params \\
\midrule
Conv-BiLSTM-Transformer & convolution, recurrence, attention & 78{,}417 \\
Conv-BiLSTM             & convolution and recurrence         & 155{,}297 \\
LSTM                    & recurrence, causal                 & 114{,}657 \\
Transformer             & self-attention                     & 101{,}793 \\
Synergy-SSM             & non-negative bottleneck with SSM   & 70{,}481 \\
Selective SSM (Mamba)   & input-dependent state space        & 30{,}770 \\
PINN                    & dynamics-constrained residual      & 40{,}132 \\
\bottomrule
\end{tabular}
\end{table}

%------------------------------------------------------------------------------
\subsection{Training Protocol}
\label{sec:training}

All models are trained with identical settings. We use the AdamW optimiser with
a base learning rate of $3 \times 10^{-4}$ and weight decay of $10^{-4}$,
minimising a Huber loss on the normalised residual. The learning rate is warmed
up linearly over the first 500 optimiser steps and thereafter controlled by a
plateau scheduler that halves the rate after five epochs without improvement in
validation loss, down to a floor of $10^{-6}$. Gradients are clipped to a maximum
norm of one. Training runs for at most 100 epochs with early stopping after 15
epochs without improvement, and the weights from the epoch with the lowest
validation loss are restored before evaluation. The batch size is 64, and the
final incomplete batch of each epoch is dropped to avoid degenerate normalisation
statistics.

Several guards protect against silent failure. Batches that produce a non-finite
loss, or whose gradient norm is not finite, are skipped rather than applied, and
training aborts if more than twenty percent of the batches in an epoch are
skipped. Training also aborts if an epoch yields no finite validation loss, and
checkpoints containing non-finite weights are refused rather than written.

Two aspects of the protocol exist specifically to make the architecture
comparison fair. First, within a fold all seven models are trained on byte
identical inputs: the base model, the calibration gains, the feature scaler and
the extracted windows are all computed once and shared. Second, the random seed
is reset to the same value before each model is constructed, so differences
between models cannot be attributed to initialisation or to data ordering. The
only variable that changes between the seven runs is the architecture itself.

%------------------------------------------------------------------------------
\subsection{Evaluation Protocol}
\label{sec:protocol}

We evaluate under two protocols that answer different questions.

\subsubsection{Within-day evaluation}

Recordings from a single session are partitioned into training, validation and
test subsets, and the model never sees data from any other session. This measures
whether the residual can be learned at all when the recording conditions are held
fixed, and provides a reference point that is free of between-session variation.
Splits are performed at the level of whole recordings. This is essential rather
than conventional: consecutive samples in these signals are correlated at
approximately 0.99, so dividing a single recording into training and test
portions would place heavily overlapping windows on both sides of the boundary
and the reported error would be optimistic. Three rotating splits are run per
session so that different recordings serve as the test set, and results are
averaged across them.

\subsubsection{Cross-day evaluation}

One session is held out entirely for testing, a second is used for validation,
and the remaining sessions are used for training. The procedure is repeated so
that each eligible session serves once as the test session. This matches the
deployment scenario, in which a model trained on previously recorded sessions is
applied to a new electrode setup.

The validation session is selected by a fixed rotation among the eligible
sessions. We deliberately do not select it by similarity to the test session.
Doing so would require knowing the test session's force statistics, which are
precisely what is unavailable when a model must be chosen for deployment, and
would make the reported performance optimistic.

One session is designated as training data only and never serves as a test or
validation session. Its gain lies a factor of roughly two below that of every
other session, so holding it out would measure extrapolation beyond the range of
gains seen in training rather than transfer between comparable sessions. Its
value lies in training, where it widens the range of conditions the model has
observed. The extrapolation case is informative in its own right and is reported
separately in Section~\ref{sec:extrapolation}.

%------------------------------------------------------------------------------
\subsection{Metrics}
\label{sec:metrics}

We report root mean squared error (RMSE), mean absolute error (MAE) and the
coefficient of determination ($R^2$), each computed per recording and then
averaged. Every metric is computed on the evaluation segment only, that is, after
the calibration segment has been excluded. For every result we report the
corresponding base model performance on exactly the same samples, so that the
contribution of the learned compensation is always visible against the physics
based baseline rather than in isolation.

RMSE is treated as the primary metric. The reason concerns the behaviour of
$R^2$ on this dataset and should be understood before reading the results. The
coefficient of determination measures error relative to the variance of the
target within each recording. Several recordings contain long low-activity
intervals and consequently have small force variance, so a modest absolute error
produces a large negative $R^2$. Averaging such values across recordings yields
aggregate figures that are negative even where absolute error is clearly reduced,
and those figures are dominated by the least informative recordings rather than
by overall accuracy. We therefore rank methods by RMSE and MAE, which are
expressed in the units of the measured force and are comparable across
recordings, and report $R^2$ alongside them for completeness rather than as the
basis for any conclusion.


%==============================================================================
%==============================================================================
% Section IV: Experiments and Results
%
% Included from main.tex via %==============================================================================
% Section IV: Experiments and Results
%
% Included from main.tex via %==============================================================================
% Section IV: Experiments and Results
%
% Included from main.tex via \input{sections/results}
% Compile main.tex, not this file.
%==============================================================================
\section{Experiments and Results}
\label{sec:results}

%------------------------------------------------------------------------------
\subsection{Experimental Design}
\label{sec:design}

Four experiments were run, summarised in Table~\ref{tab:experiments}. The first
two correspond to the evaluation protocols of Section~\ref{sec:protocol} and
form the main results. The remaining two are ablations that remove one component
of the method at a time while holding everything else fixed. All four use the
same window length, stride, optimiser settings and epoch budget, so differences
between them are attributable to the stated change alone.

\begin{table}[!t]
\caption{Experiments. All use a window of 100 samples with stride 20, a maximum
of 100 epochs with early stopping, and a \SI{10}{\second} calibration window
unless stated otherwise.}
\label{tab:experiments}
\centering
\begin{tabular}{llp{3.0cm}}
\toprule
ID & Configuration & Question addressed \\
\midrule
E1 & Cross-day, four folds & Performance in the deployment setting \\
E2 & Within-day, three folds per session & Behaviour when between-session variation is removed \\
E3 & E1 without gain calibration & Contribution of calibration \\
E4 & E1 with a static base model & Contribution of base model dynamics \\
\bottomrule
\end{tabular}
\end{table}

All experiments were run on a single NVIDIA V100 GPU. Training the seven
architectures took 122 minutes for the four cross-day folds, 88 minutes for the
fifteen within-day folds, and 128 minutes for the dynamics ablation.

%------------------------------------------------------------------------------
\subsection{Cross-Day Performance}
\label{sec:crossday}

Table~\ref{tab:crossday} reports cross-day results. Each column is a session
held out entirely from training and validation, and each entry is the mean RMSE
over the recordings of that session. The row labelled EDM base model gives the
performance of the calibrated base model on exactly the same samples, so every
compensated result can be read against its own baseline.

\begin{table*}[!t]
\caption{Cross-day (leave-one-day-out) RMSE per held-out session. Lower is
better. The best compensated result in each column is shown in bold. Every
model improves on the base model in every fold.}
\label{tab:crossday}
\centering
\begin{tabular}{lrrrrrr}
\toprule
Model & 1st\_Comp & 2nd\_Comp & 3rd\_Base & 4th\_Valid & Overall RMSE & Overall MAE \\
\midrule
EDM base model          & 0.4884 & 1.0793 & 0.1890 & 0.3659 & 0.5581 & 0.4461 \\
\midrule
Transformer             & \textbf{0.3923} & 0.9310 & \textbf{0.1342} & \textbf{0.2450} & \textbf{0.4436} & 0.3252 \\
Conv-BiLSTM             & 0.4435 & 0.9138 & 0.1436 & 0.2480 & 0.4495 & 0.3310 \\
PINN                    & 0.4384 & 0.9063 & 0.1458 & 0.2557 & 0.4501 & 0.3363 \\
LSTM                    & 0.4536 & \textbf{0.8915} & 0.1488 & 0.2638 & 0.4520 & 0.3314 \\
Conv-BiLSTM-Transformer & 0.4482 & 0.9258 & 0.1396 & 0.2521 & 0.4547 & 0.3371 \\
Selective SSM (Mamba)   & 0.4716 & 0.9010 & 0.1494 & 0.2750 & 0.4622 & 0.3448 \\
Synergy-SSM             & 0.4637 & 0.8972 & 0.1574 & 0.2830 & 0.4644 & 0.3441 \\
\bottomrule
\end{tabular}
\end{table*}

The compensator improves on the base model in every fold and for every
architecture. Taken over all held-out recordings, the best model reduces RMSE
from 0.5581 to 0.4436, a relative reduction of 20.5 percent, and even the
weakest of the seven reaches 0.4644, a reduction of 16.8 percent. Of the 224
model and recording combinations, 175 (78 percent) show lower error than the
base model on the same samples.

The size of the improvement varies with the held-out session. It is smallest for
2nd\_Comp (17.4 percent), which has the largest force range and the highest gain
of the four, and largest for 4th\_Valid (33.0 percent) and 3rd\_Base
(29.0 percent), whose gains sit near the middle of the range spanned by the
training sessions. The improvement on 1st\_Comp is 19.7 percent. This ordering is
consistent with the interpretation developed in Section~\ref{sec:extrapolation}:
the method performs best when the held-out session is well inside the range of
conditions represented in training.

A one-sided reading of these numbers would be unwarranted without testing them,
so we treat each held-out recording as a paired observation and compare each
model against the base model with a Wilcoxon signed-rank test over the 32 test
recordings. All seven architectures improve on the base model at
$p < 4 \times 10^{-4}$, with the strongest evidence for the PINN
($p = 5.9 \times 10^{-7}$, improving on 28 of 32 recordings) and the Transformer
and Conv-BiLSTM ($p = 1.3 \times 10^{-5}$, 26 of 32 each). The benefit of
compensation is therefore consistent across recordings rather than driven by a
few favourable cases.

%------------------------------------------------------------------------------
\subsection{Within-Day Performance}
\label{sec:withinday}

Table~\ref{tab:withinday} reports the within-day results, in which training and
testing use different recordings from the same session. Across all 36 held-out
recordings the picture is markedly weaker than in the cross-day setting. Only
two of the seven architectures achieve a lower aggregate RMSE than the base
model, and the best of them improves on it by 2.9 percent, compared with
20.5 percent across sessions.

\begin{table}[!t]
\caption{Within-day RMSE, averaged over three rotating folds per session. The
base model is placed in rank order so that models falling below it are visible
directly.}
\label{tab:withinday}
\centering
\begin{tabular}{lrr}
\toprule
Model & Overall RMSE & Overall MAE \\
\midrule
PINN                    & \textbf{0.4343} & 0.3354 \\
Synergy-SSM             & 0.4467 & 0.3474 \\
\textit{EDM base model} & \textit{0.4473} & \textit{0.3562} \\
Conv-BiLSTM             & 0.4481 & 0.3513 \\
Selective SSM (Mamba)   & 0.4498 & 0.3593 \\
Transformer             & 0.4610 & 0.3645 \\
LSTM                    & 0.4624 & 0.3719 \\
Conv-BiLSTM-Transformer & 0.4633 & 0.3673 \\
\bottomrule
\end{tabular}
\end{table}

This aggregate requires careful interpretation, because it is not representative
of the individual sessions. Averaging RMSE over recordings whose force ranges
differ by more than an order of magnitude gives disproportionate weight to the
session with the largest absolute errors. Here that session is 2nd\_Comp, which
contributes nine of the 36 held-out recordings and has errors near 1.1, roughly
twenty times those of 5th\_Comp. The PINN is the only architecture that improves
on the base model on 2nd\_Comp, which is the main reason it leads the aggregate.

Examined session by session, the result is more favourable. In every one of the
five sessions at least one architecture achieves a lower RMSE than the base
model: by 7.9 percent on 1st\_Comp, 3.6 percent on 2nd\_Comp, 13.2 percent on
3rd\_Base, 13.9 percent on 4th\_Valid and 11.5 percent on 5th\_Comp. What does
not hold is consistency. The architecture that performs best differs between
sessions, and no architecture improves on the base model in all five. Over all
model and recording combinations, 168 of 252 (67 percent) improve on the base
model, compared with 78 percent in the cross-day setting.

The contrast with Section~\ref{sec:crossday} is the main finding of this
experiment, and it runs opposite to the usual expectation that within-session
evaluation is the easier case. The explanation is the amount of training data
available. A cross-day fold trains on 24 to 28 recordings, whereas a within-day
fold trains on between three and nine recordings from a single session. Removing
between-session variation also removes most of the training set, and for this
task the loss of data outweighs the gain in homogeneity. A secondary factor is
that per-session calibration (Section~\ref{sec:calibration}) already absorbs much
of the between-session gain variation before the network sees the data, so the
cross-day setting is less disadvantaged than it would otherwise be.

%------------------------------------------------------------------------------
\subsection{Ablation: Per-Session Gain Calibration}
\label{sec:ablation-calib}

% PENDING. E3 did not complete: two concurrently submitted jobs wrote to the
% same output directory, and the first to finish relocated the directory while
% E3 was still running. The run must be repeated in isolation before this
% subsection can be written.
%
% Expected content once available:
%   - cross-day RMSE for base and all seven models with calibration disabled
%   - comparison against Table IV to isolate the contribution of calibration
%   - this supports the paper's central methodological claim, so it should be
%     reported with the same per-fold detail as Section IV-B

%------------------------------------------------------------------------------
\subsection{Ablation: Base Model Dynamics}
\label{sec:ablation-lpf}

Section~\ref{sec:basemodel} argued that the first-order stage is required
because the force to sEMG relationship is not single-valued. To test this, the
cross-day experiment was repeated with the first-order stage removed, leaving a
purely static saturating nonlinearity. All other settings were unchanged.

\begin{table}[!t]
\caption{Effect of removing the first-order stage from the base model. Values
are cross-day RMSE over all held-out recordings.}
\label{tab:lpf}
\centering
\begin{tabular}{lrr}
\toprule
 & With dynamics & Static only \\
\midrule
EDM base model          & 0.5581 & 0.5683 \\
Transformer             & \textbf{0.4436} & 0.4665 \\
Conv-BiLSTM             & 0.4495 & 0.4687 \\
Conv-BiLSTM-Transformer & 0.4547 & 0.4697 \\
PINN                    & 0.4501 & 0.4845 \\
LSTM                    & 0.4520 & 0.4913 \\
Synergy-SSM             & 0.4644 & 0.4917 \\
Selective SSM (Mamba)   & 0.4622 & 0.4929 \\
\bottomrule
\end{tabular}
\end{table}

Removing the dynamics degrades both stages of the method, as shown in
Table~\ref{tab:lpf}. The base model itself worsens from 0.5581 to 0.5683, and
the best compensated result worsens from 0.4436 to 0.4665, a relative increase
in error of 4.9 percent. The degradation is consistent rather than selective:
all seven architectures perform worse with the static base model, and the number
of model and recording combinations improving on the base model falls from 175
to 164 out of 224.

Two observations follow. First, the compensator does not recover what the static
base model loses, even though it has access to a 100-sample history and could in
principle represent the missing dynamics itself. Placing the dynamics in the
physics-based component is therefore more effective than leaving them to be
learned. Second, the ranking of architectures is largely preserved between the
two conditions, which suggests that the choice of architecture and the choice of
base model structure act independently.

%------------------------------------------------------------------------------
\subsection{Transfer Boundary}
\label{sec:extrapolation}

Section~\ref{sec:protocol} explained that one session is reserved for training
because its gain lies well below that of all others. Treating it as a held-out
test session instead is informative, because it shows what happens when the
method is asked to extrapolate rather than interpolate. We therefore repeated
the cross-day experiment with that session held out, keeping all other settings
identical.

The outcome reverses the main result. The calibrated base model achieves an RMSE
of 0.0663 on this session, and every one of the seven architectures performs
worse, ranging from 0.0772 for the Transformer to 0.0838 for the Synergy-SSM.
The best model is thus 16.4 percent worse than applying no compensation at all,
and only 14 of 70 model and recording combinations improve on the base model,
compared with 78 percent in Table~\ref{tab:crossday}.

The distinguishing feature of this session is that its gain of 0.44 lies a factor
of 1.9 below the lowest gain among the remaining sessions, so no training session
brackets it from below. The compensator has learned a correction appropriate to
the range of conditions it observed and applies it outside that range, where it
is systematically wrong. This delimits the applicability of the method: it
improves accuracy when the deployment session falls inside the range of
conditions spanned by training, and degrades accuracy when it does not. Because
the gain can be estimated from the calibration segment before any prediction is
made, this condition is checkable at deployment time, and falling back to the
uncompensated base model is a safe response when it is not met.

%------------------------------------------------------------------------------
\subsection{Comparison Between Architectures}
\label{sec:architecture}

The seven architectures span convolutional, recurrent, attention-based,
state-space and physics-constrained designs, and differ by a factor of five in
parameter count, from 30{,}770 to 155{,}297. Despite this, their cross-day
results occupy a narrow band: mean RMSE ranges from 0.4436 to 0.4644, a relative
spread of 4.7 percent, against a 20.5 percent improvement over the base model.
The difference between the best and worst architecture is therefore about a
quarter of the difference between using compensation and not using it.

We tested this directly. A Friedman test across the 32 held-out recordings
rejects the hypothesis that the seven architectures have identical error
distributions ($\chi^2 = 23.88$, $p = 5 \times 10^{-4}$), so the models are not
strictly interchangeable. However, when all 21 pairwise comparisons are made with
Wilcoxon signed-rank tests and corrected for multiple comparisons using the Holm
procedure, no pair differs significantly at the 5 percent level. The smallest
adjusted $p$ value is 0.221, for LSTM against the selective state-space model.

The appropriate conclusion is therefore narrower than architectural equivalence.
There is detectable heterogeneity among the seven, but the data do not identify
which architecture is preferable: no individual comparison survives correction,
and the spread is small relative to the effect of compensation itself. For this
task and dataset, the choice of sequence architecture is not the factor that
determines accuracy. Parameter count in particular is uninformative, since the
largest model ranks second and the smallest ranks sixth.

This has a practical consequence. Where two architectures cannot be separated on
accuracy, secondary criteria become decisive. The selective state-space model
attains comparable error with the fewest parameters but the longest training
time, owing to its sequential scan. The LSTM is the only strictly causal model in
the comparison and is therefore the natural candidate for a streaming
implementation. The Transformer achieves the lowest mean error and trains
quickly, but it processes the whole window bidirectionally and is less suited to
real-time use.

% Compile main.tex, not this file.
%==============================================================================
\section{Experiments and Results}
\label{sec:results}

%------------------------------------------------------------------------------
\subsection{Experimental Design}
\label{sec:design}

Four experiments were run, summarised in Table~\ref{tab:experiments}. The first
two correspond to the evaluation protocols of Section~\ref{sec:protocol} and
form the main results. The remaining two are ablations that remove one component
of the method at a time while holding everything else fixed. All four use the
same window length, stride, optimiser settings and epoch budget, so differences
between them are attributable to the stated change alone.

\begin{table}[!t]
\caption{Experiments. All use a window of 100 samples with stride 20, a maximum
of 100 epochs with early stopping, and a \SI{10}{\second} calibration window
unless stated otherwise.}
\label{tab:experiments}
\centering
\begin{tabular}{llp{3.0cm}}
\toprule
ID & Configuration & Question addressed \\
\midrule
E1 & Cross-day, four folds & Performance in the deployment setting \\
E2 & Within-day, three folds per session & Behaviour when between-session variation is removed \\
E3 & E1 without gain calibration & Contribution of calibration \\
E4 & E1 with a static base model & Contribution of base model dynamics \\
\bottomrule
\end{tabular}
\end{table}

All experiments were run on a single NVIDIA V100 GPU. Training the seven
architectures took 122 minutes for the four cross-day folds, 88 minutes for the
fifteen within-day folds, and 128 minutes for the dynamics ablation.

%------------------------------------------------------------------------------
\subsection{Cross-Day Performance}
\label{sec:crossday}

Table~\ref{tab:crossday} reports cross-day results. Each column is a session
held out entirely from training and validation, and each entry is the mean RMSE
over the recordings of that session. The row labelled EDM base model gives the
performance of the calibrated base model on exactly the same samples, so every
compensated result can be read against its own baseline.

\begin{table*}[!t]
\caption{Cross-day (leave-one-day-out) RMSE per held-out session. Lower is
better. The best compensated result in each column is shown in bold. Every
model improves on the base model in every fold.}
\label{tab:crossday}
\centering
\begin{tabular}{lrrrrrr}
\toprule
Model & 1st\_Comp & 2nd\_Comp & 3rd\_Base & 4th\_Valid & Overall RMSE & Overall MAE \\
\midrule
EDM base model          & 0.4884 & 1.0793 & 0.1890 & 0.3659 & 0.5581 & 0.4461 \\
\midrule
Transformer             & \textbf{0.3923} & 0.9310 & \textbf{0.1342} & \textbf{0.2450} & \textbf{0.4436} & 0.3252 \\
Conv-BiLSTM             & 0.4435 & 0.9138 & 0.1436 & 0.2480 & 0.4495 & 0.3310 \\
PINN                    & 0.4384 & 0.9063 & 0.1458 & 0.2557 & 0.4501 & 0.3363 \\
LSTM                    & 0.4536 & \textbf{0.8915} & 0.1488 & 0.2638 & 0.4520 & 0.3314 \\
Conv-BiLSTM-Transformer & 0.4482 & 0.9258 & 0.1396 & 0.2521 & 0.4547 & 0.3371 \\
Selective SSM (Mamba)   & 0.4716 & 0.9010 & 0.1494 & 0.2750 & 0.4622 & 0.3448 \\
Synergy-SSM             & 0.4637 & 0.8972 & 0.1574 & 0.2830 & 0.4644 & 0.3441 \\
\bottomrule
\end{tabular}
\end{table*}

The compensator improves on the base model in every fold and for every
architecture. Taken over all held-out recordings, the best model reduces RMSE
from 0.5581 to 0.4436, a relative reduction of 20.5 percent, and even the
weakest of the seven reaches 0.4644, a reduction of 16.8 percent. Of the 224
model and recording combinations, 175 (78 percent) show lower error than the
base model on the same samples.

The size of the improvement varies with the held-out session. It is smallest for
2nd\_Comp (17.4 percent), which has the largest force range and the highest gain
of the four, and largest for 4th\_Valid (33.0 percent) and 3rd\_Base
(29.0 percent), whose gains sit near the middle of the range spanned by the
training sessions. The improvement on 1st\_Comp is 19.7 percent. This ordering is
consistent with the interpretation developed in Section~\ref{sec:extrapolation}:
the method performs best when the held-out session is well inside the range of
conditions represented in training.

A one-sided reading of these numbers would be unwarranted without testing them,
so we treat each held-out recording as a paired observation and compare each
model against the base model with a Wilcoxon signed-rank test over the 32 test
recordings. All seven architectures improve on the base model at
$p < 4 \times 10^{-4}$, with the strongest evidence for the PINN
($p = 5.9 \times 10^{-7}$, improving on 28 of 32 recordings) and the Transformer
and Conv-BiLSTM ($p = 1.3 \times 10^{-5}$, 26 of 32 each). The benefit of
compensation is therefore consistent across recordings rather than driven by a
few favourable cases.

%------------------------------------------------------------------------------
\subsection{Within-Day Performance}
\label{sec:withinday}

Table~\ref{tab:withinday} reports the within-day results, in which training and
testing use different recordings from the same session. Across all 36 held-out
recordings the picture is markedly weaker than in the cross-day setting. Only
two of the seven architectures achieve a lower aggregate RMSE than the base
model, and the best of them improves on it by 2.9 percent, compared with
20.5 percent across sessions.

\begin{table}[!t]
\caption{Within-day RMSE, averaged over three rotating folds per session. The
base model is placed in rank order so that models falling below it are visible
directly.}
\label{tab:withinday}
\centering
\begin{tabular}{lrr}
\toprule
Model & Overall RMSE & Overall MAE \\
\midrule
PINN                    & \textbf{0.4343} & 0.3354 \\
Synergy-SSM             & 0.4467 & 0.3474 \\
\textit{EDM base model} & \textit{0.4473} & \textit{0.3562} \\
Conv-BiLSTM             & 0.4481 & 0.3513 \\
Selective SSM (Mamba)   & 0.4498 & 0.3593 \\
Transformer             & 0.4610 & 0.3645 \\
LSTM                    & 0.4624 & 0.3719 \\
Conv-BiLSTM-Transformer & 0.4633 & 0.3673 \\
\bottomrule
\end{tabular}
\end{table}

This aggregate requires careful interpretation, because it is not representative
of the individual sessions. Averaging RMSE over recordings whose force ranges
differ by more than an order of magnitude gives disproportionate weight to the
session with the largest absolute errors. Here that session is 2nd\_Comp, which
contributes nine of the 36 held-out recordings and has errors near 1.1, roughly
twenty times those of 5th\_Comp. The PINN is the only architecture that improves
on the base model on 2nd\_Comp, which is the main reason it leads the aggregate.

Examined session by session, the result is more favourable. In every one of the
five sessions at least one architecture achieves a lower RMSE than the base
model: by 7.9 percent on 1st\_Comp, 3.6 percent on 2nd\_Comp, 13.2 percent on
3rd\_Base, 13.9 percent on 4th\_Valid and 11.5 percent on 5th\_Comp. What does
not hold is consistency. The architecture that performs best differs between
sessions, and no architecture improves on the base model in all five. Over all
model and recording combinations, 168 of 252 (67 percent) improve on the base
model, compared with 78 percent in the cross-day setting.

The contrast with Section~\ref{sec:crossday} is the main finding of this
experiment, and it runs opposite to the usual expectation that within-session
evaluation is the easier case. The explanation is the amount of training data
available. A cross-day fold trains on 24 to 28 recordings, whereas a within-day
fold trains on between three and nine recordings from a single session. Removing
between-session variation also removes most of the training set, and for this
task the loss of data outweighs the gain in homogeneity. A secondary factor is
that per-session calibration (Section~\ref{sec:calibration}) already absorbs much
of the between-session gain variation before the network sees the data, so the
cross-day setting is less disadvantaged than it would otherwise be.

%------------------------------------------------------------------------------
\subsection{Ablation: Per-Session Gain Calibration}
\label{sec:ablation-calib}

% PENDING. E3 did not complete: two concurrently submitted jobs wrote to the
% same output directory, and the first to finish relocated the directory while
% E3 was still running. The run must be repeated in isolation before this
% subsection can be written.
%
% Expected content once available:
%   - cross-day RMSE for base and all seven models with calibration disabled
%   - comparison against Table IV to isolate the contribution of calibration
%   - this supports the paper's central methodological claim, so it should be
%     reported with the same per-fold detail as Section IV-B

%------------------------------------------------------------------------------
\subsection{Ablation: Base Model Dynamics}
\label{sec:ablation-lpf}

Section~\ref{sec:basemodel} argued that the first-order stage is required
because the force to sEMG relationship is not single-valued. To test this, the
cross-day experiment was repeated with the first-order stage removed, leaving a
purely static saturating nonlinearity. All other settings were unchanged.

\begin{table}[!t]
\caption{Effect of removing the first-order stage from the base model. Values
are cross-day RMSE over all held-out recordings.}
\label{tab:lpf}
\centering
\begin{tabular}{lrr}
\toprule
 & With dynamics & Static only \\
\midrule
EDM base model          & 0.5581 & 0.5683 \\
Transformer             & \textbf{0.4436} & 0.4665 \\
Conv-BiLSTM             & 0.4495 & 0.4687 \\
Conv-BiLSTM-Transformer & 0.4547 & 0.4697 \\
PINN                    & 0.4501 & 0.4845 \\
LSTM                    & 0.4520 & 0.4913 \\
Synergy-SSM             & 0.4644 & 0.4917 \\
Selective SSM (Mamba)   & 0.4622 & 0.4929 \\
\bottomrule
\end{tabular}
\end{table}

Removing the dynamics degrades both stages of the method, as shown in
Table~\ref{tab:lpf}. The base model itself worsens from 0.5581 to 0.5683, and
the best compensated result worsens from 0.4436 to 0.4665, a relative increase
in error of 4.9 percent. The degradation is consistent rather than selective:
all seven architectures perform worse with the static base model, and the number
of model and recording combinations improving on the base model falls from 175
to 164 out of 224.

Two observations follow. First, the compensator does not recover what the static
base model loses, even though it has access to a 100-sample history and could in
principle represent the missing dynamics itself. Placing the dynamics in the
physics-based component is therefore more effective than leaving them to be
learned. Second, the ranking of architectures is largely preserved between the
two conditions, which suggests that the choice of architecture and the choice of
base model structure act independently.

%------------------------------------------------------------------------------
\subsection{Transfer Boundary}
\label{sec:extrapolation}

Section~\ref{sec:protocol} explained that one session is reserved for training
because its gain lies well below that of all others. Treating it as a held-out
test session instead is informative, because it shows what happens when the
method is asked to extrapolate rather than interpolate. We therefore repeated
the cross-day experiment with that session held out, keeping all other settings
identical.

The outcome reverses the main result. The calibrated base model achieves an RMSE
of 0.0663 on this session, and every one of the seven architectures performs
worse, ranging from 0.0772 for the Transformer to 0.0838 for the Synergy-SSM.
The best model is thus 16.4 percent worse than applying no compensation at all,
and only 14 of 70 model and recording combinations improve on the base model,
compared with 78 percent in Table~\ref{tab:crossday}.

The distinguishing feature of this session is that its gain of 0.44 lies a factor
of 1.9 below the lowest gain among the remaining sessions, so no training session
brackets it from below. The compensator has learned a correction appropriate to
the range of conditions it observed and applies it outside that range, where it
is systematically wrong. This delimits the applicability of the method: it
improves accuracy when the deployment session falls inside the range of
conditions spanned by training, and degrades accuracy when it does not. Because
the gain can be estimated from the calibration segment before any prediction is
made, this condition is checkable at deployment time, and falling back to the
uncompensated base model is a safe response when it is not met.

%------------------------------------------------------------------------------
\subsection{Comparison Between Architectures}
\label{sec:architecture}

The seven architectures span convolutional, recurrent, attention-based,
state-space and physics-constrained designs, and differ by a factor of five in
parameter count, from 30{,}770 to 155{,}297. Despite this, their cross-day
results occupy a narrow band: mean RMSE ranges from 0.4436 to 0.4644, a relative
spread of 4.7 percent, against a 20.5 percent improvement over the base model.
The difference between the best and worst architecture is therefore about a
quarter of the difference between using compensation and not using it.

We tested this directly. A Friedman test across the 32 held-out recordings
rejects the hypothesis that the seven architectures have identical error
distributions ($\chi^2 = 23.88$, $p = 5 \times 10^{-4}$), so the models are not
strictly interchangeable. However, when all 21 pairwise comparisons are made with
Wilcoxon signed-rank tests and corrected for multiple comparisons using the Holm
procedure, no pair differs significantly at the 5 percent level. The smallest
adjusted $p$ value is 0.221, for LSTM against the selective state-space model.

The appropriate conclusion is therefore narrower than architectural equivalence.
There is detectable heterogeneity among the seven, but the data do not identify
which architecture is preferable: no individual comparison survives correction,
and the spread is small relative to the effect of compensation itself. For this
task and dataset, the choice of sequence architecture is not the factor that
determines accuracy. Parameter count in particular is uninformative, since the
largest model ranks second and the smallest ranks sixth.

This has a practical consequence. Where two architectures cannot be separated on
accuracy, secondary criteria become decisive. The selective state-space model
attains comparable error with the fewest parameters but the longest training
time, owing to its sequential scan. The LSTM is the only strictly causal model in
the comparison and is therefore the natural candidate for a streaming
implementation. The Transformer achieves the lowest mean error and trains
quickly, but it processes the whole window bidirectionally and is less suited to
real-time use.

% Compile main.tex, not this file.
%==============================================================================
\section{Experiments and Results}
\label{sec:results}

%------------------------------------------------------------------------------
\subsection{Experimental Design}
\label{sec:design}

Four experiments were run, summarised in Table~\ref{tab:experiments}. The first
two correspond to the evaluation protocols of Section~\ref{sec:protocol} and
form the main results. The remaining two are ablations that remove one component
of the method at a time while holding everything else fixed. All four use the
same window length, stride, optimiser settings and epoch budget, so differences
between them are attributable to the stated change alone.

\begin{table}[!t]
\caption{Experiments. All use a window of 100 samples with stride 20, a maximum
of 100 epochs with early stopping, and a \SI{10}{\second} calibration window
unless stated otherwise.}
\label{tab:experiments}
\centering
\begin{tabular}{llp{3.0cm}}
\toprule
ID & Configuration & Question addressed \\
\midrule
E1 & Cross-day, four folds & Performance in the deployment setting \\
E2 & Within-day, three folds per session & Behaviour when between-session variation is removed \\
E3 & E1 without gain calibration & Contribution of calibration \\
E4 & E1 with a static base model & Contribution of base model dynamics \\
\bottomrule
\end{tabular}
\end{table}

All experiments were run on a single NVIDIA V100 GPU. Training the seven
architectures took 122 minutes for the four cross-day folds, 88 minutes for the
fifteen within-day folds, and 128 minutes for the dynamics ablation.

%------------------------------------------------------------------------------
\subsection{Cross-Day Performance}
\label{sec:crossday}

Table~\ref{tab:crossday} reports cross-day results. Each column is a session
held out entirely from training and validation, and each entry is the mean RMSE
over the recordings of that session. The row labelled EDM base model gives the
performance of the calibrated base model on exactly the same samples, so every
compensated result can be read against its own baseline.

\begin{table*}[!t]
\caption{Cross-day (leave-one-day-out) RMSE per held-out session. Lower is
better. The best compensated result in each column is shown in bold. Every
model improves on the base model in every fold.}
\label{tab:crossday}
\centering
\begin{tabular}{lrrrrrr}
\toprule
Model & 1st\_Comp & 2nd\_Comp & 3rd\_Base & 4th\_Valid & Overall RMSE & Overall MAE \\
\midrule
EDM base model          & 0.4884 & 1.0793 & 0.1890 & 0.3659 & 0.5581 & 0.4461 \\
\midrule
Transformer             & \textbf{0.3923} & 0.9310 & \textbf{0.1342} & \textbf{0.2450} & \textbf{0.4436} & 0.3252 \\
Conv-BiLSTM             & 0.4435 & 0.9138 & 0.1436 & 0.2480 & 0.4495 & 0.3310 \\
PINN                    & 0.4384 & 0.9063 & 0.1458 & 0.2557 & 0.4501 & 0.3363 \\
LSTM                    & 0.4536 & \textbf{0.8915} & 0.1488 & 0.2638 & 0.4520 & 0.3314 \\
Conv-BiLSTM-Transformer & 0.4482 & 0.9258 & 0.1396 & 0.2521 & 0.4547 & 0.3371 \\
Selective SSM (Mamba)   & 0.4716 & 0.9010 & 0.1494 & 0.2750 & 0.4622 & 0.3448 \\
Synergy-SSM             & 0.4637 & 0.8972 & 0.1574 & 0.2830 & 0.4644 & 0.3441 \\
\bottomrule
\end{tabular}
\end{table*}

The compensator improves on the base model in every fold and for every
architecture. Taken over all held-out recordings, the best model reduces RMSE
from 0.5581 to 0.4436, a relative reduction of 20.5 percent, and even the
weakest of the seven reaches 0.4644, a reduction of 16.8 percent. Of the 224
model and recording combinations, 175 (78 percent) show lower error than the
base model on the same samples.

The size of the improvement varies with the held-out session. It is smallest for
2nd\_Comp (17.4 percent), which has the largest force range and the highest gain
of the four, and largest for 4th\_Valid (33.0 percent) and 3rd\_Base
(29.0 percent), whose gains sit near the middle of the range spanned by the
training sessions. The improvement on 1st\_Comp is 19.7 percent. This ordering is
consistent with the interpretation developed in Section~\ref{sec:extrapolation}:
the method performs best when the held-out session is well inside the range of
conditions represented in training.

A one-sided reading of these numbers would be unwarranted without testing them,
so we treat each held-out recording as a paired observation and compare each
model against the base model with a Wilcoxon signed-rank test over the 32 test
recordings. All seven architectures improve on the base model at
$p < 4 \times 10^{-4}$, with the strongest evidence for the PINN
($p = 5.9 \times 10^{-7}$, improving on 28 of 32 recordings) and the Transformer
and Conv-BiLSTM ($p = 1.3 \times 10^{-5}$, 26 of 32 each). The benefit of
compensation is therefore consistent across recordings rather than driven by a
few favourable cases.

%------------------------------------------------------------------------------
\subsection{Within-Day Performance}
\label{sec:withinday}

Table~\ref{tab:withinday} reports the within-day results, in which training and
testing use different recordings from the same session. Across all 36 held-out
recordings the picture is markedly weaker than in the cross-day setting. Only
two of the seven architectures achieve a lower aggregate RMSE than the base
model, and the best of them improves on it by 2.9 percent, compared with
20.5 percent across sessions.

\begin{table}[!t]
\caption{Within-day RMSE, averaged over three rotating folds per session. The
base model is placed in rank order so that models falling below it are visible
directly.}
\label{tab:withinday}
\centering
\begin{tabular}{lrr}
\toprule
Model & Overall RMSE & Overall MAE \\
\midrule
PINN                    & \textbf{0.4343} & 0.3354 \\
Synergy-SSM             & 0.4467 & 0.3474 \\
\textit{EDM base model} & \textit{0.4473} & \textit{0.3562} \\
Conv-BiLSTM             & 0.4481 & 0.3513 \\
Selective SSM (Mamba)   & 0.4498 & 0.3593 \\
Transformer             & 0.4610 & 0.3645 \\
LSTM                    & 0.4624 & 0.3719 \\
Conv-BiLSTM-Transformer & 0.4633 & 0.3673 \\
\bottomrule
\end{tabular}
\end{table}

This aggregate requires careful interpretation, because it is not representative
of the individual sessions. Averaging RMSE over recordings whose force ranges
differ by more than an order of magnitude gives disproportionate weight to the
session with the largest absolute errors. Here that session is 2nd\_Comp, which
contributes nine of the 36 held-out recordings and has errors near 1.1, roughly
twenty times those of 5th\_Comp. The PINN is the only architecture that improves
on the base model on 2nd\_Comp, which is the main reason it leads the aggregate.

Examined session by session, the result is more favourable. In every one of the
five sessions at least one architecture achieves a lower RMSE than the base
model: by 7.9 percent on 1st\_Comp, 3.6 percent on 2nd\_Comp, 13.2 percent on
3rd\_Base, 13.9 percent on 4th\_Valid and 11.5 percent on 5th\_Comp. What does
not hold is consistency. The architecture that performs best differs between
sessions, and no architecture improves on the base model in all five. Over all
model and recording combinations, 168 of 252 (67 percent) improve on the base
model, compared with 78 percent in the cross-day setting.

The contrast with Section~\ref{sec:crossday} is the main finding of this
experiment, and it runs opposite to the usual expectation that within-session
evaluation is the easier case. The explanation is the amount of training data
available. A cross-day fold trains on 24 to 28 recordings, whereas a within-day
fold trains on between three and nine recordings from a single session. Removing
between-session variation also removes most of the training set, and for this
task the loss of data outweighs the gain in homogeneity. A secondary factor is
that per-session calibration (Section~\ref{sec:calibration}) already absorbs much
of the between-session gain variation before the network sees the data, so the
cross-day setting is less disadvantaged than it would otherwise be.

%------------------------------------------------------------------------------
\subsection{Ablation: Per-Session Gain Calibration}
\label{sec:ablation-calib}

% PENDING. E3 did not complete: two concurrently submitted jobs wrote to the
% same output directory, and the first to finish relocated the directory while
% E3 was still running. The run must be repeated in isolation before this
% subsection can be written.
%
% Expected content once available:
%   - cross-day RMSE for base and all seven models with calibration disabled
%   - comparison against Table IV to isolate the contribution of calibration
%   - this supports the paper's central methodological claim, so it should be
%     reported with the same per-fold detail as Section IV-B

%------------------------------------------------------------------------------
\subsection{Ablation: Base Model Dynamics}
\label{sec:ablation-lpf}

Section~\ref{sec:basemodel} argued that the first-order stage is required
because the force to sEMG relationship is not single-valued. To test this, the
cross-day experiment was repeated with the first-order stage removed, leaving a
purely static saturating nonlinearity. All other settings were unchanged.

\begin{table}[!t]
\caption{Effect of removing the first-order stage from the base model. Values
are cross-day RMSE over all held-out recordings.}
\label{tab:lpf}
\centering
\begin{tabular}{lrr}
\toprule
 & With dynamics & Static only \\
\midrule
EDM base model          & 0.5581 & 0.5683 \\
Transformer             & \textbf{0.4436} & 0.4665 \\
Conv-BiLSTM             & 0.4495 & 0.4687 \\
Conv-BiLSTM-Transformer & 0.4547 & 0.4697 \\
PINN                    & 0.4501 & 0.4845 \\
LSTM                    & 0.4520 & 0.4913 \\
Synergy-SSM             & 0.4644 & 0.4917 \\
Selective SSM (Mamba)   & 0.4622 & 0.4929 \\
\bottomrule
\end{tabular}
\end{table}

Removing the dynamics degrades both stages of the method, as shown in
Table~\ref{tab:lpf}. The base model itself worsens from 0.5581 to 0.5683, and
the best compensated result worsens from 0.4436 to 0.4665, a relative increase
in error of 4.9 percent. The degradation is consistent rather than selective:
all seven architectures perform worse with the static base model, and the number
of model and recording combinations improving on the base model falls from 175
to 164 out of 224.

Two observations follow. First, the compensator does not recover what the static
base model loses, even though it has access to a 100-sample history and could in
principle represent the missing dynamics itself. Placing the dynamics in the
physics-based component is therefore more effective than leaving them to be
learned. Second, the ranking of architectures is largely preserved between the
two conditions, which suggests that the choice of architecture and the choice of
base model structure act independently.

%------------------------------------------------------------------------------
\subsection{Transfer Boundary}
\label{sec:extrapolation}

Section~\ref{sec:protocol} explained that one session is reserved for training
because its gain lies well below that of all others. Treating it as a held-out
test session instead is informative, because it shows what happens when the
method is asked to extrapolate rather than interpolate. We therefore repeated
the cross-day experiment with that session held out, keeping all other settings
identical.

The outcome reverses the main result. The calibrated base model achieves an RMSE
of 0.0663 on this session, and every one of the seven architectures performs
worse, ranging from 0.0772 for the Transformer to 0.0838 for the Synergy-SSM.
The best model is thus 16.4 percent worse than applying no compensation at all,
and only 14 of 70 model and recording combinations improve on the base model,
compared with 78 percent in Table~\ref{tab:crossday}.

The distinguishing feature of this session is that its gain of 0.44 lies a factor
of 1.9 below the lowest gain among the remaining sessions, so no training session
brackets it from below. The compensator has learned a correction appropriate to
the range of conditions it observed and applies it outside that range, where it
is systematically wrong. This delimits the applicability of the method: it
improves accuracy when the deployment session falls inside the range of
conditions spanned by training, and degrades accuracy when it does not. Because
the gain can be estimated from the calibration segment before any prediction is
made, this condition is checkable at deployment time, and falling back to the
uncompensated base model is a safe response when it is not met.

%------------------------------------------------------------------------------
\subsection{Comparison Between Architectures}
\label{sec:architecture}

The seven architectures span convolutional, recurrent, attention-based,
state-space and physics-constrained designs, and differ by a factor of five in
parameter count, from 30{,}770 to 155{,}297. Despite this, their cross-day
results occupy a narrow band: mean RMSE ranges from 0.4436 to 0.4644, a relative
spread of 4.7 percent, against a 20.5 percent improvement over the base model.
The difference between the best and worst architecture is therefore about a
quarter of the difference between using compensation and not using it.

We tested this directly. A Friedman test across the 32 held-out recordings
rejects the hypothesis that the seven architectures have identical error
distributions ($\chi^2 = 23.88$, $p = 5 \times 10^{-4}$), so the models are not
strictly interchangeable. However, when all 21 pairwise comparisons are made with
Wilcoxon signed-rank tests and corrected for multiple comparisons using the Holm
procedure, no pair differs significantly at the 5 percent level. The smallest
adjusted $p$ value is 0.221, for LSTM against the selective state-space model.

The appropriate conclusion is therefore narrower than architectural equivalence.
There is detectable heterogeneity among the seven, but the data do not identify
which architecture is preferable: no individual comparison survives correction,
and the spread is small relative to the effect of compensation itself. For this
task and dataset, the choice of sequence architecture is not the factor that
determines accuracy. Parameter count in particular is uninformative, since the
largest model ranks second and the smallest ranks sixth.

This has a practical consequence. Where two architectures cannot be separated on
accuracy, secondary criteria become decisive. The selective state-space model
attains comparable error with the fewest parameters but the longest training
time, owing to its sequential scan. The LSTM is the only strictly causal model in
the comparison and is therefore the natural candidate for a streaming
implementation. The Transformer achieves the lowest mean error and trains
quickly, but it processes the whole window bidirectionally and is less suited to
real-time use.

%==============================================================================
\section{Discussion}
\label{sec:discussion}
% - Calibration, not architecture, is what makes cross-day compensation work.
% - Training volume dominates domain match (cross-day beats within-day).
% - The unexplained residual sits well above the measurement noise floor
%   (9-23% of residual RMSE) yet is poorly predicted from sEMG alone,
%   suggesting missing inputs: joint angle and contraction velocity enter force
%   production through force-length and force-velocity relations, and the base
%   model has no such dependence.
% - Practical implication: deployment protocols should include a brief
%   reference contraction when electrodes are fitted.

%==============================================================================
\section{Limitations}
\label{sec:limitations}
% - Single sEMG channel; limits the synergy formulation (Section III-H5).
% - Single subject.
% - The ~13x monotonic gain drift across sessions is unexplained; whether the
%   force transducer was recalibrated between sessions could not be verified.
% - Selective SSM implemented without fused kernels (runtime only).
% - Single random seed; seed variance not characterised.

%==============================================================================
\section{Conclusion}
\label{sec:conclusion}

%==============================================================================
\section*{Acknowledgment}
% Computation performed on the National Computational Infrastructure (NCI)
% Gadi supercomputer, project mi23.

\bibliographystyle{IEEEtran}
\bibliography{references}

\end{document}
